% !TeX root = ../../Book3.tex
\chapter{\texorpdfstring{\(\Gamma\)}{\AnalysisPdfGamma} 函数与 Mellin 方法}
\label{b3:ch:21}
\BookChapterNumber{b3:ch:21}
\begin{chapterabstract}
Gamma 函数把阶乘递推、参数积分与整函数乘积联系起来。本章从 Euler 积分的两端控制出发，证明全纯性、递推和严格对数凸性，用 Bohr--Mollerup 定理说明正实轴上的唯一性，再通过递推和 Taylor 扣除完成亚纯延拓。Euler 极限确定倒数乘积的规范化，进而给出无零点性、对数导数、反射与乘法公式；正实轴 Stirling 公式则从积分峰值附近的尺度得到增长常数。最后转入 Mellin 变换，借对数坐标明确权重、Fourier 归一化与反演条件，并以 Gamma 的留数展开说明端点渐近如何成为变换的极点。
\end{chapterabstract}

\begin{learninggoals}
\begin{itemize}
\item 分别控制 Euler 积分在零与无穷远的行为，验证参数微分，并说明递推、正性、规范化与对数凸性各自承担的作用。
\item 用递推粘合亚纯延拓，计算极点留数及有限部分，区分收敛积分、Taylor 正则化与延拓后的函数值。
\item 从 Euler 极限得到倒数乘积，先证明非零性再取对数导数，核对 Euler 常数和局部幂和系数的来源。
\item 保留 Beta 换元的 Jacobi 因子以及反射、倍增、乘法公式的规范化常数，区分正实轴渐近和特殊竖线估计。
\item 从两端可积性确定 Mellin 收敛带，用对数坐标解释缩放、权重、乘法卷积与 Fourier 变换之间的关系。
\item 分开验证绝对反演、有限移线和余项极限的条件，追踪幂次、对数因子、极点位置与极点阶数的对应。
\end{itemize}
\end{learninggoals}

若只要求 \(F(n+1)=n!\) 和 \(F(z+1)=zF(z)\)，阶乘在复平面中的延拓并不唯一：乘上适当的一周期函数，这些条件仍然成立。Euler 积分提供了一个自然候选，但要理解它为何成为通常的 Gamma 函数，还须辨明积分带来的正性、凸性和增长信息。另一方面，同一个积分中的幂 \(t^{z-1}\) 也能用来分析一般函数；把 \(e^{-t}\) 换成 \(f(t)\)，就得到本章后半部的 Mellin 变换。

\ProseParagraphBreak
\secref{b3:sec:2101}先处理原积分。原点附近由实部的正下界控制，无穷远由指数下降控制；加入对数幂后仍可积，才允许逐次对参数求导。递推来自边界项消失的分部积分，Gaussian 积分确定半整数处的常数。阅读 Bohr--Mollerup 证明时，应注意凸性比较的是哪些割线，以及整数点已知值怎样夹住非整数处的值。固定实平移的 Gamma 比值可以直接从这些凸性估计得到，不必先使用 Stirling 公式。

\ProseParagraphBreak
\secref{b3:sec:2102}把递推转为向左延拓的工具。不同移动次数必须在重叠区域给出同一个函数；分母零点处则要检查分子是否也消失，才能确认极点阶数。Taylor 扣除提供另一种具体表示：先减去造成发散的低次项，再把它们的积分写成明确的有理奇部。这个过程解释了延拓值从何而来，却不使原先发散的 Euler 积分自动收敛。

\ProseParagraphBreak
\secref{b3:sec:2103}从有限积分推得 Euler 极限，再整理出 \(1/\Gamma\) 的正规乘积。一次指数修正保证尾项可和，Euler 常数记录有限乘积规范化的极限。这里先比较两个有限表达式及其极限，证明 \(\Gamma\) 无零点，然后才使用对数导数。由此得到的 \(\psi_0\)、高阶导数与局部幂和，也能回到参数积分中解释对数矩和曲率。

\ProseParagraphBreak
\secref{b3:sec:2104}将这些表示互相连接。相反参数的乘积给反射公式，Beta 积分的二维换元给倍增关系，递推与对数凸性则固定一般乘法公式中的函数自由度。Stirling 证明另从积分的峰值出发：在宽度为 \(\sqrt{x}\) 的中部出现 Gaussian 极限，两侧尾部需要各自估计。该证明只建立正实轴版本；整数和半整数竖线上的模值由反射与递推独立求出，不以未证明的复参数渐近代替。

\ProseParagraphBreak
\secref{b3:sec:2105}沿每条实部为 \(c\) 的竖线，把 Mellin 变换写成 \(e^{cu}f(e^u)\) 的 Fourier 变换。改变竖线也就改变权重，因此收敛带属于变换资料的一部分。绝对反演要求竖线上的变换可积；围道移动还要控制水平边；若继续让新竖线趋向左方，又需一条随移动次数变化的余项界。指数函数的例子将这三步分别完成，最后再把端点的幂与对数展开解释为亚纯变换的主部。

\ProseParagraphBreak
本章使用参数积分、支配收敛、全纯函数的局部一致极限与留数定理，以及\chapref{b3:ch:20}的正规乘积和正弦乘积。Mellin 反演所需的 Fourier 归一化会在使用处重新写明。Stein、Shakarchi 的《Complex Analysis》\cite{Stein2003Complex}可供继续阅读特殊函数与后续数论问题的联系；本章所用的 Gamma 恒等式、反演条件和移线估计均在正文中建立。

% !TeX root = ../../Book3.tex
\section{Euler \texorpdfstring{\(\Gamma\)}{\AnalysisPdfGamma} 函数}
\label{b3:sec:2101}

阶乘满足乘法递推。若要把整数变量换成实变量或复变量，需要先找到一个在整数处给出阶乘、同时又能对参数作分析运算的表达。Euler 积分提供了这样的函数，但递推和整数点取值并不足以唯一确定它。本节先从积分得到全纯性、递推和半整数值，再用正实轴上的对数凸性建立唯一性与平移比值估计。

\subsection{Euler 积分}
\label{b3:subsec:210101}

积分 \(\int_0^\infty e^{-t}t^n\,\mathrm{d}t\) 中，分部积分会把幂次降低一，并产生因子 \(n\)。为了使递推写成变量本身乘以函数，定义时把幂次取为 \(z-1\)。

\Needspace{6\baselineskip}
\begin{definition}\label{b3:def:21-gamma}
对 \(z\in\mathbb C\)、\(\operatorname{Re}z>0\)，定义\term[gamma-hanshu]{Gamma 函数}[gamma function]
\[
 \Gamma(z)\coloneqq\int_0^\infty e^{-t}t^{z-1}\,\mathrm{d}t,
 \quad t^{z-1}=\exp\bigl((z-1)\log t\bigr),\quad t>0,
\]
其中 \(\log t\) 是正实轴上的实对数。
\end{definition}
\symidx{\(\Gamma(z)\)：Gamma 函数}
积分变量只取正实数，所以复幂在这里没有绕过原点选择分支的问题。若 \(z=\sigma+\mathrm{i}y\)，则 \(|t^{z-1}|=t^{\sigma-1}\)，虚部只产生模为一的振荡因子。

\ProseParagraphBreak

两端的收敛机制不同。在 \((0,1)\) 上，\(e^{-1}\le e^{-t}\le1\)，所以绝对可积恰好要求 \(\sigma>0\)。在 \([1,\infty)\) 上，指数衰减压过任意固定实次幂，对所有实数 \(\sigma\) 都可积。因此 \(\operatorname{Re}z>0\) 是这个积分的绝对收敛域。要在积分号下对 \(z\) 求导，还须控制每次求导带来的 \(\log t\)。

\Needspace{6\baselineskip}
\begin{proposition}[全纯性与参数导数]\label{b3:prop:21-euler-derivatives}
Euler 积分定义的 \(\Gamma\) 在半平面 \(\operatorname{Re}z>0\) 全纯。对每个非负整数 \(k\)，有
\[
 \Gamma^{(k)}(z)
 =\int_0^\infty e^{-t}t^{z-1}(\log t)^k\,\mathrm{d}t,
 \quad \operatorname{Re}z>0,
\]
且这些积分在每个紧参数集上绝对一致收敛。此外，\(\Gamma(x)>0\) 对 \(x>0\) 成立，并且
\[
 |\Gamma(\sigma+\mathrm{i}y)|\le\Gamma(\sigma),
 \quad \sigma>0,\quad y\in\mathbb R.
\]
\end{proposition}
\begin{proof}
固定紧集 \(K\subset\{\operatorname{Re}z>0\}\)，选 \(0<a\le\operatorname{Re}z\le b\) 对所有 \(z\in K\) 成立。在 \(0<t<1\) 上，
\[
 |e^{-t}t^{z-1}(\log t)^k|
 \le t^{a-1}|\log t|^k.
\]
代换 \(u=-\log t\)，再作有限次分部积分，得到
\[
 \int_0^1t^{a-1}|\log t|^k\,\mathrm{d}t
 =\int_0^\infty e^{-au}u^k\,\mathrm{d}u
 =\frac{k!}{a^{k+1}}<\infty.
\]
在 \(t\ge1\) 上，则有
\[
 |e^{-t}t^{z-1}(\log t)^k|
 \le e^{-t}t^{b-1}(\log t)^k
 \le C_{b,k}e^{-t/2}.
\]
最后一个常数存在，因为 \(t^{b-1}(\log t)^k e^{-t/2}\) 在 \([1,\infty)\) 有界。这些控制函数不依赖 \(z\in K\)。

截断积分 \(G_N(z)=\int_{1/N}^{N}e^{-t}t^{z-1}\,\mathrm{d}t\) 是整函数，在有限积分区间可以逐次对参数求导。上述共同控制表明，\(G_N\) 及其各个固定阶导数的积分表达在 \(K\) 上一致收敛。由\cref{b3:prop:04-derivative-convergence}，极限全纯，导数正是陈述中的积分。正实参数时，被积函数处处为正，故 \(\Gamma(x)>0\)。最后，对 Euler 积分使用三角不等式即得竖线上的模估计。
\end{proof}
这里的一致性针对紧参数集；当实部趋近零时，共同控制中的 \(a\) 也会趋近零，不能保留同一个可积上界。右半平面上的竖线估计只给出有界性，尚未给出随 \(|y|\) 增大的衰减速度。

\subsection{函数方程}
\label{b3:subsec:210102}

积分定义使变量平移一步恰好对应一次分部积分。端点项是否消失，仍由原点的幂次和无穷远的指数分别控制。

\Needspace{6\baselineskip}
\begin{proposition}[递推与归一化]\label{b3:prop:21-recurrence}
对所有 \(z\in\mathbb C\)、\(\operatorname{Re}z>0\)，有
\[
 \Gamma(z+1)=z\Gamma(z),\quad \Gamma(1)=1.
\]
因此，对每个非负整数 \(n\)，\(\Gamma(n+1)=n!\)。
\end{proposition}
\begin{proof}
在 \([\varepsilon,R]\) 上分部积分，得到
\[
 \int_\varepsilon^R e^{-t}t^z\,\mathrm{d}t
 =\bigl[-e^{-t}t^z\bigr]_\varepsilon^R
  +z\int_\varepsilon^R e^{-t}t^{z-1}\,\mathrm{d}t.
\]
设 \(\sigma=\operatorname{Re}z>0\)。无穷远端的边界项模为 \(e^{-R}R^\sigma\to0\)，原点端的模为 \(e^{-\varepsilon}\varepsilon^\sigma\to0\)。两个积分都绝对收敛，所以可分别令 \(R\to\infty\)、\(\varepsilon\downarrow0\)，得到递推式。直接计算 \(\int_0^\infty e^{-t}\,\mathrm{d}t=1\) 给出归一化，再作整数次递推得到阶乘。
\end{proof}
归一化选定了乘法常数，却没有排除一周期因子。例如，固定 \(0<\eta<1\)，函数
\[
 F(z)=\Gamma(z)\bigl(1+\eta\sin(2\pi z)\bigr)
\]
在右半平面全纯，满足同一递推、同一整数点取值，而且在正实轴仍为正。它与 \(\Gamma\) 不同，因此即使同时要求全纯性和正实轴上的正性，递推仍不能给出唯一性。

\ProseParagraphBreak

递推还把半整数值归约到一个起始值。这个值由 Gaussian 积分确定，适合在讨论渐近公式之前独立算出。

\Needspace{6\baselineskip}
\begin{proposition}[Gaussian 积分与半整数值]\label{b3:prop:21-half-integers}
对每个非负整数 \(n\)，有
\[
 \Gamma\left(\frac12\right)=\sqrt\pi,
 \quad
 \Gamma\left(n+\frac12\right)
 =\frac{(2n)!}{4^n n!}\sqrt\pi.
\]
\end{proposition}
\begin{proof}
令 \(I=\int_{\mathbb R}e^{-u^2}\,\mathrm{d}u\)。这个积分有限且为正：在 \(|u|\ge1\) 上可由 \(e^{-|u|}\) 控制，在有限区间被积函数有界。被积函数非负，Tonelli 定理给
\[
 I^2=\int_{\mathbb R^2}e^{-(u^2+v^2)}\,\mathrm{d}u\,\mathrm{d}v.
\]
在圆环上使用极坐标换元，Jacobi 因子为 \(r\)，再让内半径趋零、外半径趋于无穷。非负性允许用单调收敛取极限，于是
\[
 I^2=\int_0^{2\pi}\int_0^\infty e^{-r^2}r\,\mathrm{d}r\,\mathrm{d}\theta
 =2\pi\cdot\frac12=\pi.
\]
由 \(I>0\) 确定 \(I=\sqrt\pi\)。在 Euler 积分中令 \(t=u^2\)，得到
\[
 \Gamma\left(\frac12\right)
 =2\int_0^\infty e^{-u^2}\,\mathrm{d}u=I.
\]
最后，重复使用递推式，
\[
 \Gamma\left(n+\frac12\right)
 =\left(\prod_{j=0}^{n-1}\left(j+\frac12\right)\right)\sqrt\pi
 =\frac{1\cdot3\cdots(2n-1)}{2^n}\sqrt\pi
 =\frac{(2n)!}{4^n n!}\sqrt\pi.
\]
\(n=0\) 时有限乘积为空，公式仍成立。
\end{proof}
同一换元给出 Gaussian 偶数矩
\[
 \int_{\mathbb R}u^{2n}e^{-u^2}\,\mathrm{d}u
 =\Gamma\left(n+\frac12\right),\quad n\ge0.
\]
奇数矩因被积函数为奇函数而为零；这些积分均绝对收敛。阶乘和半整数值由此分别对应指数权重与 Gaussian 权重下的整数矩。

\subsection{对数凸性}
\label{b3:subsec:210103}

Euler 积分中的参数出现在指数 \(t^{x-1}\) 上，因此参数的凸组合会变成两个正被积函数的几何平均。Hölder 不等式正好比较这种几何平均的积分。

\Needspace{6\baselineskip}
\begin{proposition}[严格对数凸性]\label{b3:prop:21-logconvex}
Euler 积分定义的实函数 \(x\mapsto\log\Gamma(x)\) 在 \((0,\infty)\) 上严格凸。即对 \(x,y>0\)、\(x\ne y\) 和 \(0<\lambda<1\)，
\[
 \Gamma(\lambda x+(1-\lambda)y)
 <\Gamma(x)^\lambda\Gamma(y)^{1-\lambda}.
\]
\end{proposition}
\begin{proof}
令 \(A(t)=e^{-t}t^{x-1}\)、\(B(t)=e^{-t}t^{y-1}\)。二者严格为正且可积，并且
\[
 e^{-t}t^{\lambda x+(1-\lambda)y-1}
 =A(t)^\lambda B(t)^{1-\lambda}.
\]
对右边使用指数 \(1/\lambda\) 与 \(1/(1-\lambda)\) 的 Hölder 不等式，得
\[
 \int_0^\infty A^\lambda B^{1-\lambda}\,\mathrm{d}t
 \le\left(\int_0^\infty A\,\mathrm{d}t\right)^\lambda
      \left(\int_0^\infty B\,\mathrm{d}t\right)^{1-\lambda}.
\]
在两个非零正函数之间，等号要求 \(A\) 与 \(B\) 几乎处处成比例，即 \(t^{x-y}\) 几乎处处为常数。但 \(x\ne y\) 时该函数在 \((0,\infty)\) 严格单调，不可能满足这个条件。因此不等式严格，取对数即得结论。
\end{proof}
还可以把严格性读成一个二阶导数公式。固定 \(x>0\)，以
\[
 w_x(t)=\frac{e^{-t}t^{x-1}}{\Gamma(x)},\quad t>0,
\]
作为概率密度，并令 \(m_x=\int_0^\infty(\log t)w_x(t)\,\mathrm{d}t\)。参数导数公式保证这些对数矩有限，且
\[
 \frac{\mathrm{d}^2}{\mathrm{d}x^2}\log\Gamma(x)
 =\frac{\Gamma''(x)}{\Gamma(x)}
  -\left(\frac{\Gamma'(x)}{\Gamma(x)}\right)^2
 =\int_0^\infty(\log t-m_x)^2w_x(t)\,\mathrm{d}t>0.
\]
最后一项是 \(\log t\) 的方差；密度在整个正半轴严格为正，而 \(\log t\) 不恒定，所以方差严格为正。归一化的分母 \(\Gamma(x)\) 是这一步所必需的，单独观察 \(\Gamma''(x)>0\) 只能得到普通凸性。

\subsection{Bohr–Mollerup 定理}
\label{b3:subsec:210104}

递推固定了整数点的对数差值，而凸性控制整数之间的割线斜率。把一个未知的 \(F(x)\) 向右递推到 \(F(n+x)\)，再利用相邻整数之间的斜率夹逼，就能让未知量落入相对宽度趋于零的区间。这给出\term[bohr-mollerup]{Bohr--Mollerup 定理}[Bohr--Mollerup theorem]。

\Needspace{6\baselineskip}
\begin{theorem}[Bohr--Mollerup]\label{b3:thm:21-bohr-mollerup}
设 \(F:(0,\infty)\to(0,\infty)\) 满足
\[
 F(1)=1,\quad F(x+1)=xF(x)\quad(x>0),
\]
并且 \(x\mapsto\log F(x)\) 凸。则对每个 \(x>0\)，\(F(x)=\Gamma(x)\)。
\end{theorem}
\begin{proof}
令 \(h=\log F\)。递推与归一化先给 \(F(n)=(n-1)!\)，因而对整数 \(n\ge2\)，
\[
 h(n)-h(n-1)=\log(n-1),\quad
 h(n+1)-h(n)=\log n.
\]
固定 \(0<x\le1\)。凸函数的割线斜率随区间向右移动而非减，所以
\[
 h(n)-h(n-1)
 \le\frac{h(n+x)-h(n)}x
 \le h(n+1)-h(n).
\]
右侧亦可由 \(n+x=(1-x)n+x(n+1)\) 的凸性直接得到。于是
\[
 x\log(n-1)\le h(n+x)-h(n)\le x\log n.
\]
指数化后，再用 \(F(n+x)=F(x)x(x+1)\cdots(x+n-1)\)，可得
\[
 A_n(x)\le F(x)\le B_n(x),
\]
其中
\[
 A_n(x)=\frac{(n-1)!(n-1)^x}{x(x+1)\cdots(x+n-1)},\quad
 B_n(x)=\frac{(n-1)!n^x}{x(x+1)\cdots(x+n-1)}.
\]
这两个端点都严格为正，且 \(B_n(x)/A_n(x)=(n/(n-1))^x\to1\)。

已经由积分证明的 \(\Gamma\) 同样满足归一化、递推和对数凸性，所以也落在同一个区间 \([A_n(x),B_n(x)]\) 中。两次夹逼相除，得到
\[
 \left(\frac{n-1}{n}\right)^x
 \le\frac{F(x)}{\Gamma(x)}
 \le\left(\frac n{n-1}\right)^x.
\]
令 \(n\to\infty\)，即得 \(F(x)=\Gamma(x)\) 于 \((0,1]\)。任意正实数可写成 \(x+m\)，其中 \(0<x\le1\)、\(m\ge0\) 为整数，再用有限次递推便扩展到全部正实轴。
\end{proof}
证明只要求 \(\log F\) 凸，不预先要求 \(F\) 可微。若另外给定一个在右半平面全纯的候选函数，且它在正实轴满足定理的条件，则先由本定理识别实轴上的值，再由恒等定理识别整个右半平面内的函数。实变量唯一性与复变量唯一性分别由这两步承担。

\ProseParagraphBreak

对数凸性也能给出精确到相对误差的平移比值。这里的自变量始终沿正实轴趋于无穷，平移参数则保持固定。

\Needspace{6\baselineskip}
\begin{proposition}[固定实平移的比值]\label{b3:prop:21-shift-ratio}
设 \(x>0\)、\(0\le a\le1\)。则
\[
 \left(\frac{x}{x+a}\right)^{1-a}
 \le\frac{\Gamma(x+a)}{x^a\Gamma(x)}\le1.
\]
特别地，对任意固定的 \(a,b\in\mathbb R\)，当 \(x\to+\infty\) 时有
\[
 \frac{\Gamma(x+a)}{\Gamma(x+b)}
 =x^{a-b}\bigl(1+O_{a,b}(x^{-1})\bigr).
\]
后一公式在所有参数 \(x+a,x+b\) 均为正的充分大半轴上使用。
\end{proposition}
\begin{proof}
\(a=0,1\) 时由归一化的比值和递推直接成立。对 \(0<a<1\)，将 \(x+a=(1-a)x+a(x+1)\) 代入对数凸性，得到
\[
 \Gamma(x+a)\le\Gamma(x)^{1-a}\Gamma(x+1)^a
 =x^a\Gamma(x).
\]
另一方面，\(x+1=a(x+a)+(1-a)(x+a+1)\)，故
\[
 \Gamma(x+1)
 \le\Gamma(x+a)^a\Gamma(x+a+1)^{1-a}
 =(x+a)^{1-a}\Gamma(x+a).
\]
使用 \(\Gamma(x+1)=x\Gamma(x)\) 并整理，便得下界。固定 \(a\in[0,1]\) 时，下界为 \((1+a/x)^{-(1-a)}=1+O_a(x^{-1})\)，所以
\[
 \frac{\Gamma(x+a)}{\Gamma(x)}=x^a\bigl(1+O_a(x^{-1})\bigr).
\]

任意固定实数 \(c\) 可写成 \(c=k+\theta\)，其中 \(k\in\mathbb Z\)、\(0\le\theta<1\)。若 \(k\ge0\)，有限次递推给
\[
 \Gamma(x+c)=\Gamma(x+\theta)
             \prod_{j=0}^{k-1}(x+\theta+j).
\]
若 \(k=-q<0\)，则
\[
 \Gamma(x+c)=\frac{\Gamma(x+\theta)}
 {\prod_{j=0}^{q-1}(x+c+j)}.
\]
对充分大的 \(x\)，这些因子均为正；每个因子提出一个 \(x\) 后，有限乘积或其倒数均为 \(1+O_c(x^{-1})\)。结合已经证明的 \(\theta\) 情形，得到 \(\Gamma(x+c)/\Gamma(x)=x^c(1+O_c(x^{-1}))\)。分别取 \(c=a,b\)，再相除，即得一般比值。
\end{proof}
误差常数依赖固定的平移参数；这条结论不能直接用于随 \(x\) 增大的平移量。它也没有给出复变量在任意方向趋于无穷的比值渐近。

\ProseParagraphBreak

半整数值与比值估计合起来，已经能够确定中心二项式系数的主项。对正整数 \(n\)，
\[
 \binom{2n}{n}
 =\frac{4^n}{\sqrt\pi}
   \frac{\Gamma(n+1/2)}{\Gamma(n+1)}
 =\frac{4^n}{\sqrt{\pi n}}\bigl(1+O(n^{-1})\bigr).
\]
这里常数 \(\sqrt\pi\) 来自 Gaussian 积分，比值的幂次来自对数凸性与递推。两种信息共同确定了这个离散系数的渐近尺度。

% !TeX root = ../../Book3.tex
\section{亚纯延拓}
\label{b3:sec:2102}

递推式既能向右计算，也能向左延拓。后一用法必须检查不同移动次数给出的值相同，并辨明分母的零点是真极点还是可去点。

\subsection{递推关系}
\label{b3:subsec:210201}

对任意正整数 \(N\)，在 \(\operatorname{Re}z>-N\) 定义
\[
\Gamma(z)=\frac{\Gamma(z+N)}{z(z+1)\cdots(z+N-1)}.
\]
当分母不为零时这是全纯函数，其他点至多为极点。用一次递推消去 \(z+N\) 可知，采用 \(N+1\) 与 \(N\) 的定义在重叠区域一致。因此这些局部表达式粘合成整个复平面上的亚纯函数，并继续满足递推式。以后 \(\Gamma\) 均指这一延拓。

\subsection{极点}
\label{b3:subsec:210202}

\begin{proposition}\label{b3:prop:21-poles}
Gamma 函数的极点恰为 \(0,-1,-2,\ldots\)，每个都是一阶极点。
\end{proposition}
\begin{proof}
若 \(z\) 不是非正整数，选足够大的 \(N\) 可使分子处于右半平面，且分母各因子非零，故该点全纯。在 \(z=-m\) 附近取 \(N=m+1\)。分子在该点等于 \(\Gamma(1)=1\)，分母只有因子 \(z+m\) 为零，其余因子都非零，所以是恰好一阶的极点。
\end{proof}
延拓并不赋予原积分在 \(\operatorname{Re}z\le0\) 上收敛的含义。例如 \(\Gamma(-1/2)=-2\sqrt\pi\)，而形式积分 \(\int_0^\infty e^{-t}t^{-3/2}\,\mathrm{d}t\) 发散且被积函数为正；两者不能混为一谈。

\subsection{留数}
\label{b3:subsec:210203}

上述局部表示直接给出
\[
\operatorname{Res}_{z=-m}\Gamma(z)=\frac{(-1)^m}{m!},\quad m=0,1,2,\ldots.
\]
事实上，提出 \(z+m\) 后，其他因子在 \(-m\) 的乘积为 \((-1)^m m!\)。在反演积分中移动围道时，这串留数将产生幂级数的系数，而极点位置给出相应的幂次。

\ProseParagraphBreak

复共轭关系 \(\Gamma(\overline z)=\overline{\Gamma(z)}\) 先由右半平面的积分得到，再经恒等定理延拓。负实轴相邻极点之间的符号因此可用递推明确计算：在 \((-m-1,-m)\) 上符号为 \((-1)^{m+1}\)。此符号变化来自越过一阶极点，不意味着存在实零点；事实上下一节将证明复零点也不存在。

\ProseParagraphBreak

极点附近还能得到有限部分，而不仅是留数。先由乘积的首项展开得到 \(\Gamma(z)=z^{-1}-\gamma+O(z)\)。在 \(z=-m+\varepsilon\) 使用递推，分母非零部分为
\[
 \prod_{j=0}^{m-1}(-m+\varepsilon+j)
 =(-1)^m m!\left(1-H_m\varepsilon+O(\varepsilon^2)\right).
\]
分子 \(\Gamma(1+\varepsilon)=1-\gamma\varepsilon+O(\varepsilon^2)\)。因此
\[
 \Gamma(-m+\varepsilon)=\frac{(-1)^m}{m!}
       \left(\frac1\varepsilon+H_m-\gamma+O(\varepsilon)\right),
 \qquad H_0=0.
\]
当把 Gamma 与另一个含零点的函数相乘时，留数决定可去后的值，有限部分则决定导数。只知道极点阶数不足以完成这些局部计算。
\ProseParagraphBreak
参数正则化提供一个与递推不同的延拓图像。在 \(-1<\operatorname{Re}z<0\) 中，积分
\[
 \int_0^1(e^{-t}-1)t^{z-1}\,\mathrm{d}t
 +\int_1^\infty e^{-t}t^{z-1}\,\mathrm{d}t+\frac1z
\]
绝对收敛并等于 Gamma 延拓。因为 \(e^{-t}-1=O(t)\)，第一项比原 Euler 积分多出一个可积幂次；最后的 \(1/z\) 恰把减掉的常数项补回。继续减去更多 Taylor 项，就能向更左的半平面推进。所谓解析延拓在这里表现为一项明确的奇部加一个普通收敛积分，并没有把发散积分本身赋予任意的数值。

% !TeX root = ../../Book3.tex
\section{乘积表示}
\label{b3:sec:2103}

上一节确定了全部极点，但还没有排除复零点。有限乘积的极限既给出零点信息，也解释 Euler 常数为什么出现在 Gamma 函数中。

\subsection{Weierstrass 乘积}
\label{b3:subsec:210301}

\begin{lemma}[Euler 极限]\label{b3:lem:21-euler-limit}
在 \(\operatorname{Re}z>0\) 上局部一致地有
\[
\Gamma(z)=\lim_{N\to\infty}\frac{N!N^z}{z(z+1)\cdots(z+N)}.
\]
\end{lemma}
\begin{proof}
连续分部积分给
\(\int_0^1 t^{z-1}(1-t)^N\,\mathrm{d}t=N!/[z(z+1)\cdots(z+N)]\)。乘以 \(N^z\)，换元 \(u=Nt\)，所得积分为
\(\int_0^N u^{z-1}(1-u/N)^N\,\mathrm{d}u\)。在 \(0<u<N\) 上有 \((1-u/N)^N\le e^{-u}\)，而它逐点趋于 \(e^{-u}\)。对紧集 \(a\le\operatorname{Re}z\le b\)，用 \((u^{a-1}+u^{b-1})e^{-u}\) 控制。将积分分成固定紧区间与共同小尾部，即得对 \(z\) 的局部一致收敛。
\end{proof}
这个证明同时保留复参数与一致性，避免只在实数处建立公式后未经说明就逐项微分。有限分部积分使用的端点条件来自 \(\operatorname{Re}z>0\)。

\subsection{Euler 常数}
\label{b3:subsec:210302}

记 \(H_N=\sum_{n=1}^N1/n\)。序列 \(H_N-\log N\) 单调递减，因为
\(1/(N+1)\le\log(1+1/N)\)，并由积分比较有下界零，故极限存在，记为\term[euler-changshu]{Euler 常数}[Euler constant] \(\gamma\)。把有限乘积写为
\[
\frac{z(z+1)\cdots(z+N)}{N!N^z}
=z\exp\bigl(z(H_N-\log N)\bigr)
\prod_{n=1}^N(1+z/n)e^{-z/n}.
\]
紧集上，尾部对数 \(\log(1+z/n)-z/n=O_K(n^{-2})\) 一致可和；前面有限个因子无需取对数。于是右侧在全平面局部一致收敛为整函数。

\subsection{\texorpdfstring{\(1/\Gamma\)}{1/\AnalysisPdfGamma} 作为整函数}
\label{b3:subsec:210303}

\begin{theorem}[Gamma 函数的 Weierstrass 乘积]\label{b3:thm:21-product}
Gamma 函数没有零点，并且
\[
\frac1{\Gamma(z)}=ze^{\gamma z}\prod_{n=1}^\infty(1+z/n)e^{-z/n}.
\]
右端整函数的零点恰为非正整数，且均为单零点。
\end{theorem}
\begin{proof}
设右端为 \(P(z)\)。Euler 极限与相应的有限倒数乘积在右半平面都局部一致收敛，它们的乘积恒为一，取极限得 \(P\Gamma=1\)，不必事先假设 \(\Gamma\ne0\)。尾部指数的和在避开有限个头部零点的邻域全纯，因此尾乘积不为零；全部零点正是所列的线性因子零点，重数为一。亚纯恒等定理把 \(P\Gamma=1\) 延拓到全平面，给出结论。
\end{proof}
对数导数在避开极点的紧集可逐项求导。若记 \(\psi_0(z)=\Gamma'(z)/\Gamma(z)\)，则
\[
\psi_0(z)=-\gamma-\frac1z+\sum_{n=1}^\infty\left(\frac1n-\frac1{n+z}\right).
\]
这里用下标避免与后文的素数计数函数 \(\psi(x)\) 混淆。代入 \(z=1\) 得 \(\Gamma'(1)=-\gamma\)，再用递推得 \(\psi_0(n)=H_{n-1}-\gamma\)。整数点的导数遂由同一个调和和描述。

\ProseParagraphBreak

在零点附近展开对数乘积，可以得到 Gamma 的局部 Taylor 系数。对 \(|z|<1\)，\(\Gamma(1+z)\) 无零点且在0点等于一，取在那里为零的对数。利用 \(\Gamma(1+z)=z\Gamma(z)\)，有
\[
 \log\Gamma(1+z)=-\gamma z+
   \sum_{k=2}^\infty\frac{(-1)^k}{k}
       \left(\sum_{n=1}^\infty\frac1{n^k}\right)z^k.
\]
证明中的换序来自 \(|z|\le r<1\) 时
\(\sum_n\sum_{k\ge2}|z/n|^k/k\le C_r\sum_n n^{-2}\)。内层数列以后记为 \(\zeta(k)\)，当前只需把它看成收敛的正项数列。常数项与一次项已经分别由 \(\Gamma(1)=1\) 和 Euler 常数决定，二次项才开始包含这些幂和。
\ProseParagraphBreak
特别地，\(\Gamma''(1)=\gamma^2+\sum n^{-2}\)。与积分微分公式比较，就得到
\[
 \int_0^\infty e^{-t}\log t\,\mathrm{d}t=-\gamma,
 \qquad
 \int_0^\infty e^{-t}(\log t)^2\,\mathrm{d}t
 =\gamma^2+\sum_{n\ge1}n^{-2}.
\]
第二式减去第一式的平方，正好是前面概率解释中 \(x=1\) 的方差。乘积的系数、参数积分的矩和凸性的二阶导数在同一恒等式中相遇。
\ProseParagraphBreak
对正实 \(x\)，对数导数还给出
\[
 \psi_0'(x)=\sum_{n=0}^\infty\frac1{(n+x)^2},\qquad
 \psi_0^{(j)}(x)=(-1)^{j+1}j!\sum_{n=0}^\infty(n+x)^{-j-1}\quad(j\ge1).
\]
这些级数在任何 \(x\ge a>0\) 的紧区间上一致可微，因而 \(\psi_0'\) 正，后续导数交替符号。对数凸性的严格性由此再次可见。用递减函数的积分比较还有
\(1/x\le\psi_0'(x)\le1/x+1/x^2\)，它量化了大参数时曲率逐渐减小的速度。
\ProseParagraphBreak
有限乘积中的指数修正也有必要性。若直接写 \(\prod_{n\ge1}(1+z/n)\)，对正实 \(z\) 其对数含发散的 \(z\sum1/n\)，因而不会趋于有限非零值。乘 \(e^{-z/n}\) 消去这一不可和的一次项，剩下二次误差才可和；前面的 \(e^{\gamma z}\) 则记录有限乘积正规化的极限。把该因子删去会保留零点，却破坏与 Euler 积分的对应及递推常数。

% !TeX root = ../../Book3.tex
\section{经典公式}
\label{b3:sec:2104}

乘积能把相反参数联系起来；凸性则能把相差有理数的参数联系起来。增长估计另需考察积分质量集中的位置，不能只由函数方程猜测。

\subsection{反射公式}
\label{b3:subsec:210401}

\begin{theorem}[Euler 反射公式]\label{b3:thm:21-reflection}
对 \(z\notin\mathbb Z\)，有 \(\Gamma(z)\Gamma(1-z)=\pi/\sin\pi z\)。
\end{theorem}
\begin{proof}
对 \(1/\Gamma(z)\) 与 \(1/\Gamma(-z)\) 的正规乘积相乘，指数因子相消，得
\[
\frac1{\Gamma(z)\Gamma(-z)}=-z^2\prod_{n\ge1}(1-z^2/n^2)
=-\frac{z\sin\pi z}{\pi}.
\]
最后一步使用\secref{b3:sec:2005}已证的正弦乘积。由 \(\Gamma(1-z)=-z\Gamma(-z)\) 即得公式。
\end{proof}
这称为\term[euler-fanshe-gongshi]{Euler 反射公式}[Euler reflection formula]。取 \(z=1/2\)，积分的正性确定 \(\Gamma(1/2)=\sqrt\pi\)，排除平方根符号的不确定性。

\subsection{倍增公式}
\label{b3:subsec:210402}

\begin{theorem}[Legendre 倍增公式]\label{b3:thm:21-duplication}
作为亚纯恒等式，有
\[
\Gamma(z)\Gamma(z+1/2)=2^{1-2z}\sqrt\pi\,\Gamma(2z).
\]
\end{theorem}
\begin{proof}
对 \(x>0\) 定义 \(F(x)=2^{x-1}\Gamma(x/2)\Gamma((x+1)/2)/\sqrt\pi\)。直接递推得 \(F(x+1)=xF(x)\)，且 \(F(1)=1\)。函数 \(\log F\) 为两个凸函数与线性函数之和，因此凸。Bohr--Mollerup 定理给 \(F(x)=\Gamma(x)\)；令 \(x=2z\)，先在正实轴成立，再由亚纯恒等定理扩展。
\end{proof}
该式称为\term[legendre-beizeng-gongshi]{Legendre 倍增公式}[Legendre duplication formula]。例如它把半整数阶乘写成
\(\Gamma(n+1/2)=(2n)!\sqrt\pi/(4^nn!)\)，便于精确计算高斯积分的偶数矩。

\subsection{乘法公式}
\label{b3:subsec:210403}

对整数 \(m\ge1\)，\term[gauss-chengfa-gongshi]{Gauss 乘法公式}[Gauss multiplication formula]为
\begin{equation}\label{b3:eq:21-multiplication}
\prod_{k=0}^{m-1}\Gamma(z+k/m)
=(2\pi)^{(m-1)/2}m^{1/2-mz}\Gamma(mz).
\end{equation}
先用与倍增公式相同的方法，令
\(F(x)=C m^x\prod_{k=0}^{m-1}\Gamma((x+k)/m)\)。相邻乘积约去后得到 \(F(x+1)=xF(x)\)，且对数凸，选 \(C\) 使 \(F(1)=1\) 即知它等于 \(\Gamma(x)\)。求常数时，将反射公式对 \(k=1,\ldots,m-1\) 相乘，使用
\[
\prod_{k=1}^{m-1}\sin(\pi k/m)=m/2^{m-1}.
\]
这个有限恒等式来自 \((X^m-1)/(X-1)\) 在 \(X=1\) 的值：各非平凡单位根与一的距离为 \(2\sin(\pi k/m)\)，取模即可。故
\(\prod_{k=1}^{m-1}\Gamma(k/m)=(2\pi)^{(m-1)/2}/\sqrt m\)，平方根取正值。代回 \(C\)，得到所示常数。

\subsection{Stirling 公式}
\label{b3:subsec:210404}

\begin{theorem}[Stirling]\label{b3:thm:21-stirling}
当实数 \(x\to+\infty\) 时，
\[
\Gamma(x)=\sqrt{2\pi}\,x^{x-1/2}e^{-x}(1+o(1)).
\]
\end{theorem}
\begin{proof}
在 \(\Gamma(x+1)\) 的积分中令 \(t=x+u\sqrt x\)，提出 \(x^xe^{-x}\sqrt x\)。剩余积分为
\[
I_x=\int_{-\sqrt x}^{\infty}\exp\left[x\left(\log(1+u/\sqrt x)-u/\sqrt x\right)\right]\,\mathrm{d}u.
\]
固定 \(u\) 时指数趋于 \(-u^2/2\)。对 \(|v|\le1/2\)，Taylor 定理及二阶导数给 \(\log(1+v)-v\le-cv^2\)，故中部由 \(e^{-cu^2}\) 控制。设 \(h(v)=v-\log(1+v)\)。对 \(v\ge1/2\)，有 \(h'(v)=v/(1+v)\ge1/3\)，于是右尾积分不超过
\(3x^{-1/2}e^{-xh(1/2)}\)。对 \(-1<v\le-1/2\)，\(h'(v)\le-1\)，左尾同样不超过 \(x^{-1/2}e^{-xh(-1/2)}\)。先用支配收敛处理扩大的中部，再用这两个指数小界删去尾部，得 \(I_x\to\int_{\mathbb R}e^{-u^2/2}\,\mathrm{d}u=\sqrt{2\pi}\)。最后除以 \(x\) 即得结论。
\end{proof}
这给出\term[stirling-gongshi]{Stirling 公式}[Stirling formula]的正实轴版本。证明解释了平方根因子：积分主要来自峰值 \(t=x\) 周围宽度为 \(\sqrt x\) 的区间。它不声称在穿过负实轴的复扇形上一致成立；复参数渐近展开还必须指定对数分支及与负实轴的角距离。NIST 的《Digital Library of Mathematical Functions》\cite[Sections 5.5 and 5.11]{NationalInstituteofStandardsNIST}分别列出上述函数关系和复参数渐近版本，使用时应保留其区域条件。


\ProseParagraphBreak

反射公式能给虚方向上的精确模值，而不需要复 Stirling 公式。对实 \(t\ne0\)，利用共轭关系与递推，有
\[
 |\Gamma(\mathrm{i}t)|^2=\frac{\pi}{t\sinh\pi t},\qquad
 |\Gamma(1+\mathrm{i}t)|^2=\frac{\pi t}{\sinh\pi t}.
\]
确实，反射式在 \(z=\mathrm{i}t\) 给分母 \(\sin\pi\mathrm{i}t=\mathrm{i}\sinh\pi t\)，而 \(\Gamma(1-\mathrm{i}t)=-\mathrm{i}t\Gamma(-\mathrm{i}t)\)；移项后得到第一式，第二式再乘 \(t^2\)。当 \(t\to0\)，第二式趋于一，恰与 \(\Gamma(1)=1\) 相符；当 \(|t|\to\infty\)，它给 \(|\Gamma(1+\mathrm{i}t)|\sim\sqrt{2\pi|t|}e^{-\pi|t|/2}\)。这是一条竖线上完整的衰减结论，不被误写为全扇形的统一渐近。
\ProseParagraphBreak
在半整数竖线上，同理可直接将 \(z=1/2+\mathrm{i}t\) 代入反射公式，得到
\[
 |\Gamma(1/2+\mathrm{i}t)|^2=\frac\pi{\cosh\pi t}.
\]
再对实部作整数递推，便能得到所有整数或半整数竖线上的精确乘积表达式。这些特殊线已经足以检验很多 Mellin 反演中的绝对收敛。若需要任意实部的统一估计，则必须另作插值或扇形分析，当前几条特殊线的公式并不自动包含这一推广。

% !TeX root = ../../Book3.tex
\section{Mellin 变换}
\label{b3:sec:2105}

平移问题自然使用 Fourier 变换，伸缩问题则更适合先取对数坐标。正半轴上的乘法 \(x\mapsto ax\) 在 \(u=\log x\) 中变成平移 \(u\mapsto u+\log a\)；测度 \(\mathrm dx/x\) 也恰好变成 \(\mathrm du\)。Mellin 变换由这个换元产生，它的复参数实部选择权重，虚部记录对数变量的振荡。两者承担不同作用，反演时必须同时保留。

\subsection{定义}
\label{b3:subsec:210501}

\Needspace{6\baselineskip}
\begin{definition}[Mellin 变换与绝对收敛带]\label{b3:def:21-mellin}
设 \(f:(0,\infty)\to\mathbb C\) 可测。对每个使积分绝对收敛的 \(s\in\mathbb C\)，定义其\term[mellin-bianhuan]{Mellin 变换}[Mellin transform]为
\[
 \mathcal Mf(s)\coloneqq\int_0^\infty f(x)x^{s-1}\,\mathrm dx,
 \quad x^{s-1}=\exp((s-1)\log x),
\]
其中 \(\log x\) 为正实轴上的实对数。对实数 \(c\)，记
\[
 \|f\|_{1,c}\coloneqq\int_0^\infty|f(x)|x^{c-1}\,\mathrm dx,
 \quad L^1_c=\{\,f:\|f\|_{1,c}<\infty\,\},
\]
函数按几乎处处相等识别。若集合 \(I_f=\{\,c\in\mathbb R:\|f\|_{1,c}<\infty\,\}\) 的内部为非空区间 \((a,b)\)，则称 \(a<\operatorname{Re}s<b\) 为 \(f\) 的绝对收敛带；允许 \(a=-\infty\) 或 \(b=+\infty\)。
\end{definition}
\symidx{\(\mathcal Mf(s)\)：Mellin 变换}
\symidx{\(\lVert f\rVert_{1,c}\)：Mellin 权重的绝对可积范数}
绝对值满足 \(|x^{s-1}|=x^{\operatorname{Re}s-1}\)，所以是否绝对收敛只取决于实部。若在一个实部收敛，在同一条竖线的每一点都收敛；这与函数在一个开区域全纯是两件不同的事。

\Needspace{6\baselineskip}
\begin{proposition}[收敛区间与带内全纯性]\label{b3:prop:21-mellin-strip}
设 \(f:(0,\infty)\to\mathbb C\) 可测。其绝对可积权重集合 \(I_f\) 是一个区间，允许为空、单点及具有任意端点开闭性。若实数 \(a<b\) 满足 \(f\in L^1_a\cap L^1_b\)，则 \(\mathcal Mf\) 在闭带 \(a\le\operatorname{Re}s\le b\) 有界且连续，在开带内全纯，并且对每个非负整数 \(k\)，
\[
 (\mathcal Mf)^{(k)}(s)
 =\int_0^\infty f(x)x^{s-1}(\log x)^k\,\mathrm dx,
 \quad a<\operatorname{Re}s<b.
\]
这些导数积分在开带的紧集上绝对一致收敛。
\end{proposition}
\begin{proof}
若 \(a\le\sigma\le b\)，则在 \(0<x\le1\) 上 \(x^{\sigma-1}\le x^{a-1}\)，在 \(x\ge1\) 上 \(x^{\sigma-1}\le x^{b-1}\)。因此
\[
 |f(x)x^{s-1}|
 \le |f(x)|\bigl(x^{a-1}\cdot\mathbf1_{(0,1]}(x)
                  +x^{b-1}\cdot\mathbf1_{[1,\infty)}(x)\bigr).
\]
右边可积且与 \(s\) 无关。这既证明中间权重均可积，也由支配收敛给出闭带连续性及上界 \(|\mathcal Mf(s)|\le\|f\|_{1,a}+\|f\|_{1,b}\)。所以 \(I_f\) 具有区间性质。

在开带内固定一条较小闭带 \(a+\delta\le\operatorname{Re}s\le b-\delta\)，其中 \(0<2\delta<b-a\)。对任意固定 \(k\)，函数 \(x^\delta|\log x|^k\) 在 \((0,1]\) 有界，\(x^{-\delta}(\log x)^k\) 在 \([1,\infty)\) 有界。因此加入 \((\log x)^k\) 后，仍由上述两个端点权重的常数倍控制。截断到 \([\epsilon,R]\) 的积分是 \(s\) 的整函数，其各阶导数恰由插入 \((\log x)^k\) 得到。可积控制使这些积分及导数在较小闭带上一致收敛，故可取全纯极限并逐次微分。
\end{proof}

端点不应从带内结论猜测。给定实数 \(c_0\)，函数
\[
 f(x)=\frac{x^{-c_0}}{(1+|\log x|)^2}
\]
满足 \(I_f=\{c_0\}\)：令 \(u=\log x\)，其加权绝对积分为
\[
 \int_{\mathbb R}\frac{e^{(c-c_0)u}}{(1+|u|)^2}\,\mathrm du,
\]
仅在 \(c=c_0\) 时有限。它在一条竖线上有变换，却没有非空的绝对收敛开带。相反，若 \(f\) 可积且支集包含于某个 \([\delta,R]\subset(0,\infty)\)，则所有实权重都可积，变换为整函数。端点行为决定参数域，不能只看被积函数中出现的复幂。

\subsection{缩放}
\label{b3:subsec:210502}

Mellin 变换的换元规则应连同权重一起记忆。正半轴上的伸缩不改变复参数，乘上一个幂则移动参数所在的直线；取倒数坐标会交换原点与无穷远。

\Needspace{6\baselineskip}
\begin{proposition}[缩放、幂乘与倒数换元]\label{b3:prop:21-mellin-rules}
设 \(f:(0,\infty)\to\mathbb C\) 可测，\(a>0\)、\(\lambda\in\mathbb C\)。凡右侧积分绝对收敛，下列左侧积分也绝对收敛，且有
\[
 \begin{aligned}
 \mathcal M[f(a\,\cdot)](s)&=a^{-s}\mathcal Mf(s),\\
 \mathcal M[x^\lambda f(x)](s)&=\mathcal Mf(s+\lambda),\\
 \mathcal M[f(1/x)](s)&=\mathcal Mf(-s).
 \end{aligned}
\]
全部正实数的复幂均用实对数定义。另设实数 \(c\) 和函数 \(f\in C^1((0,\infty))\) 满足
\[
 f,\;xf'\in L^1_c,
 \quad \lim_{x\downarrow0}x^cf(x)=\lim_{x\to\infty}x^cf(x)=0.
\]
则在 \(\operatorname{Re}s=c\) 上有
\begin{equation}\label{b3:eq:21-mellin-euler-operator}
 \mathcal M[xf'](s)=-s\mathcal Mf(s).
\end{equation}
\end{proposition}
\begin{proof}
第一式令 \(y=ax\)，有 \(x^{s-1}\,\mathrm dx=a^{-s}y^{s-1}\,\mathrm dy\)；对绝对值作同一换元也验证收敛性。第二式直接合并指数，绝对权重由 \(\operatorname{Re}s\) 移到 \(\operatorname{Re}s+\operatorname{Re}\lambda\)。第三式令 \(y=1/x\)，反转端点后得到 \(x^{s-1}\,\mathrm dx=-y^{-s-1}\,\mathrm dy\)，正是参数 \(-s\) 的积分。

最后在 \([\epsilon,R]\) 分部积分：
\[
 \int_\epsilon^R f'(x)x^s\,\mathrm dx
 =[x^sf(x)]_\epsilon^R-s\int_\epsilon^R f(x)x^{s-1}\,\mathrm dx.
\]
因 \(|x^{s-c}|=1\)，端点假设使边界项趋零；两项绝对可积假设允许令 \(\epsilon\downarrow0\)、\(R\to\infty\)，得到\eqref{b3:eq:21-mellin-euler-operator}。
\end{proof}

算子 \(x\,\mathrm d/\mathrm dx\) 是伸缩在参数原点的导数，因为 \(\left.\mathrm d f(e^u x)/\mathrm du\right|_{u=0}=xf'(x)\)。它对应乘子 \(-s\)。重复使用此式时，每一级函数都必须满足相应加权可积性和边界条件；形式地把微分方程换成代数方程，并不会自动消去端点贡献。

\subsection{与 Gamma 函数的关系}
\label{b3:subsec:210503}

Euler 积分就是最基本的 Mellin 变换对：
\[
 \mathcal M[e^{-x}](s)=\Gamma(s),\quad \operatorname{Re}s>0.
\]
原点处 \(e^{-x}\) 趋于一，要求正实部；无穷远的指数衰减则容许任意固定幂权重。再令 \(t=ax^p\)，得到
\begin{equation}\label{b3:eq:21-mellin-exponential}
 \mathcal M[e^{-ax^p}](s)=\frac1p a^{-s/p}\Gamma(s/p),
 \quad a,p>0,\quad \operatorname{Re}s>0.
\end{equation}
其中 \(1/p\) 与 \(a^{-s/p}\) 都由微分元和幂权重共同产生，不能只改变 Gamma 的参数而漏掉常数。

\ProseParagraphBreak

一个同时检验两端的例子是 \(f(x)=(1+x)^{-a}\)，其中 \(a>0\) 为实数。原点要求 \(\operatorname{Re}s>0\)，无穷远要求 \(\operatorname{Re}s<a\)。在这条带内作 \(u=x/(1+x)\)，由\cref{b3:prop:21-beta-integral}，
\begin{equation}\label{b3:eq:21-mellin-beta}
 \int_0^\infty\frac{x^{s-1}}{(1+x)^a}\,\mathrm dx
 =B(s,a-s)=\frac{\Gamma(s)\Gamma(a-s)}{\Gamma(a)}.
\end{equation}
取 \(a=1\)，反射公式将右边化成 \(\pi/\sin\pi s\)。它延拓为亚纯函数，但原积分依然只在 \(0<\operatorname{Re}s<1\) 绝对收敛；例如在实数 \(s>1\) 时，原积分于无穷远发散，不等于延拓的有限值。

\ProseParagraphBreak

一般的端点扣除也可在此统一说明。局部幂项 \(x^\alpha\) 在 \((0,1)\) 上给 \(1/(s+\alpha)\)，它的对数幂则给更高阶极点。为把这项计算变成延拓结论，必须同时控制扣除后的余项与未改变的无穷远积分。

\Needspace{6\baselineskip}
\begin{proposition}[端点幂与对数项的 Mellin 延拓]\label{b3:prop:21-mellin-endpoint}
设 \(f:(0,\infty)\to\mathbb C\) 可测且局部可积。设 \(A,b\in\mathbb R\)，有限个复数 \(\alpha_j\) 满足 \(\operatorname{Re}\alpha_j<A\)，并给定非负整数 \(q_j\) 及复数 \(c_{j,\ell}\)。假设当 \(x\downarrow0\) 时，
\[
 f(x)=\sum_{j=1}^J\sum_{\ell=0}^{q_j}
            c_{j,\ell}x^{\alpha_j}(\log x)^\ell+O(x^A),
\]
且对每个实数 \(\sigma<b\)，有
\[
 \int_1^\infty|f(x)|x^{\sigma-1}\,\mathrm dx<\infty.
\]
当 \(J\ge1\) 时，假设 \(\max_j(-\operatorname{Re}\alpha_j)<b\)；当 \(J=0\) 时，假设 \(-A<b\)。则原变换在 \(J\ge1\) 时于 \(\max_j(-\operatorname{Re}\alpha_j)<\operatorname{Re}s<b\) 绝对收敛，并在 \(-A<\operatorname{Re}s<b\) 有亚纯延拓
\begin{equation}\label{b3:eq:21-mellin-endpoint-extension}
 \mathcal Mf(s)=H(s)+\sum_{j=1}^J\sum_{\ell=0}^{q_j}
       \frac{(-1)^\ell\ell!c_{j,\ell}}{(s+\alpha_j)^{\ell+1}},
\end{equation}
其中 \(H\) 在这条较大开带全纯。没有主项时，原变换已经在该较大开带全纯。列出的极点只表示可能的极点；相同位置的系数须先合并，可能发生抵消。
\end{proposition}
\begin{proof}
在 \((0,1)\) 上令
\[
 r(x)=f(x)-\sum_{j=1}^J\sum_{\ell=0}^{q_j}
                   c_{j,\ell}x^{\alpha_j}(\log x)^\ell.
\]
原点附近有 \(|r(x)|\le Cx^A\)，而在远离零的紧区间，\(r\) 仍可积。所以
\[
 H(s)=\int_0^1r(x)x^{s-1}\,\mathrm dx
      +\int_1^\infty f(x)x^{s-1}\,\mathrm dx
\]
在 \(-A<\operatorname{Re}s<b\) 绝对收敛。每条内闭带都能在左右各留少量幂次余量，从而吸收任意固定次 \(|\log x|\)；与\cref{b3:prop:21-mellin-strip}相同的支配论证给出全纯性。

在原来各主项都可积的开带内，对 \(\operatorname{Re}(s+\alpha_j)>0\)，先积分 \(x^{s+\alpha_j-1}\)，再在参数上微分 \(\ell\) 次，得到
\[
 \int_0^1x^{s+\alpha_j-1}(\log x)^\ell\,\mathrm dx
 =\frac{(-1)^\ell\ell!}{(s+\alpha_j)^{\ell+1}}.
\]
参数微分仍由少量幂次余量控制。将有限个主项补回，便在非空的原收敛带得到\eqref{b3:eq:21-mellin-endpoint-extension}，其右端遂给出唯一的亚纯延拓。
\end{proof}

对 \(e^{-x}\) 扣除 \(N\) 个 Taylor 项，\(N\ge1\)，就是
\[
 \Gamma(s)=H_N(s)+\sum_{k=0}^{N-1}\frac{(-1)^k}{k!(s+k)},
 \quad \operatorname{Re}s>-N,
\]
其中 \(H_N\) 全纯。这从 Mellin 角度再次解释了 Gamma 的极点与留数：幂次 \(k\) 产生位置 \(-k\)，Taylor 系数 \((-1)^k/k!\) 变成留数。该例在无穷远有任意幂权重的可积性，因此可将命题分别用于任意大的有限 \(b\)。

\ProseParagraphBreak

这里只证明了“有受控端点展开，因此有这些极点”的方向。反向从极点读取函数渐近时，还须有可用的反演和围道移动估计。一个亚纯表达式本身没有记载积分在高处的增长，不能仅列出留数就宣布已恢复原函数。

\subsection{Mellin 反演}
\label{b3:subsec:210504}

固定一条绝对收敛竖线 \(\operatorname{Re}s=c\)。令
\[
 g_c(u)=e^{cu}f(e^u),\quad u\in\mathbb R.
\]
对数换元给出精确的范数关系与变换关系：
\begin{equation}\label{b3:eq:21-mellin-fourier}
 \begin{aligned}
 \|g_c\|_{L^1(\mathbb R)}&=\|f\|_{1,c},\\
 \mathcal Mf(c+\mathrm it)&=\int_{\mathbb R}g_c(u)e^{\mathrm itu}\,\mathrm du
                         =\widehat g_c(-t/(2\pi)).
 \end{aligned}
\end{equation}
这里 \(\widehat g(\xi)=\int g(u)e^{-2\pi\mathrm iu\xi}\,\mathrm du\) 是本书的 Fourier 归一化。正号 \(e^{\mathrm itu}\) 和因子 \(2\pi\) 决定反演中的符号与常数。

\Needspace{6\baselineskip}
\begin{theorem}[Mellin 反演]\label{b3:thm:21-mellin-inversion}
设 \(f:(0,\infty)\to\mathbb C\) 可测，\(c\in\mathbb R\)，且
\[
 \int_0^\infty|f(x)|x^{c-1}\,\mathrm dx<\infty,
 \quad \int_{\mathbb R}|\mathcal Mf(c+\mathrm it)|\,\mathrm dt<\infty,
\]
其中 \(\mathcal Mf(s)=\int_0^\infty f(x)e^{(s-1)\log x}\,\mathrm dx\)。则
\begin{equation}\label{b3:eq:21-mellin-inversion}
 f(x)=\frac1{2\pi}\int_{\mathbb R}
               \mathcal Mf(c+\mathrm it)x^{-c-\mathrm it}\,\mathrm dt
\end{equation}
对几乎处处的 \(x>0\) 成立，右端定义 \((0,\infty)\) 上的连续函数。若 \(f\) 连续，等式对每个 \(x>0\) 成立。
\end{theorem}
\begin{proof}
令 \(g(u)=e^{cu}f(e^u)\) 以及 \(F(t)=\mathcal Mf(c+\mathrm it)\)。由\eqref{b3:eq:21-mellin-fourier}，\(g\in L^1(\mathbb R)\)，且 \(F(t)=\int g(v)e^{\mathrm itv}\,\mathrm dv\)。对 \(\epsilon>0\)，在反演积分中暂时插入 Gaussian 因子，定义
\[
 J_\epsilon(u)=\frac1{2\pi}\int_{\mathbb R}
                   F(t)e^{-\epsilon t^2}e^{-\mathrm itu}\,\mathrm dt.
\]
代入 \(F\) 后，绝对二重积分由
\[
 \|g\|_1\int_{\mathbb R}e^{-\epsilon t^2}\,\mathrm dt<\infty
\]
控制，故 Fubini 定理适用。

为核对内积分的常数，记 \(A_\epsilon(w)=\int e^{-\epsilon t^2}e^{-\mathrm itw}\,\mathrm dt\)。参数微分及一次分部积分给
\[
 A_\epsilon'(w)=-\frac{w}{2\epsilon}A_\epsilon(w),
 \quad A_\epsilon(0)=\sqrt{\pi/\epsilon}.
\]
分部积分的端点项因 Gaussian 衰减而消失；参数微分由 \(|t|e^{-\epsilon t^2}\) 控制。解这个实参数微分方程，得到
\[
 \frac1{2\pi}A_\epsilon(w)
 =K_\epsilon(w)\coloneqq\frac{e^{-w^2/(4\epsilon)}}{\sqrt{4\pi\epsilon}},
 \quad J_\epsilon=K_\epsilon*g.
\]
核 \(K_\epsilon\) 非负、积分为一，是固定的径向非增可积核在尺度 \(\sqrt\epsilon\) 下的缩放。第一卷命题14.4.6说明：这样的核作用于全空间 \(L^1\) 函数，在每个 Lebesgue 点收敛到函数值。这里 \(g\in L^1(\mathbb R)\)，核的全部条件也已核对，所以
\[
 J_\epsilon(u)\longrightarrow g(u)
\]
在 \(g\) 的几乎所有点成立。采用无限支集的 Gaussian 核时，全空间可积性仍不可省去；只知道核的质量向原点集中，不足以直接套用这一逐点结论。

另一方面，假设 \(F\in L^1\) 允许以 \(|F(t)|\) 作支配函数，得到对每个 \(u\) 都成立的极限
\[
 J_\epsilon(u)\longrightarrow
 J(u)=\frac1{2\pi}\int_{\mathbb R}F(t)e^{-\mathrm itu}\,\mathrm dt.
\]
同一个支配收敛论证说明 \(J\) 连续。于是 \(g=J\) 几乎处处；令 \(u=\log x\)，再乘 \(x^{-c}\)，便得\eqref{b3:eq:21-mellin-inversion}。这一换元保持零测集结论，因为指数映射在每个有界区间 Lipschitz，而实轴可被可数个有界区间覆盖。若 \(f\) 连续，两个连续函数几乎处处相等就必处处相等。
\end{proof}

反演也可写为
\[
 f(x)=\frac1{2\pi\mathrm i}
       \int_{c-\mathrm i\infty}^{c+\mathrm i\infty}\mathcal Mf(s)x^{-s}\,\mathrm ds.
\]
直线由下到上，\(\mathrm ds=\mathrm i\,\mathrm dt\)，所以两种记法的常数一致。它是绝对收敛的竖线积分；尚未允许改变直线的位置。

\Needspace{6\baselineskip}
\begin{corollary}[一条竖线上的唯一性]\label{b3:cor:21-mellin-uniqueness}
设 \(c\in\mathbb R\)，可测函数 \(f,h\in L^1_c\)。若
\[
 \mathcal Mf(c+\mathrm it)=\mathcal Mh(c+\mathrm it),\quad t\in\mathbb R,
\]
则 \(f=h\) 几乎处处；若二者连续，则处处相等。
\end{corollary}
\begin{proof}
差 \(f-h\) 属于 \(L^1_c\)，其竖线变换恒为零，因而自动属于 \(L^1(\mathbb R,\mathrm dt)\)。对差应用反演定理，就得到零函数。这里不需要分别假设两个原变换在竖线上可积。
\end{proof}

竖线上的绝对可积性不是加权可积性的自动后果。例如 \(f=\mathbf1_{(0,1)}\) 在每个 \(c>0\) 属于 \(L^1_c\)，且 \(\mathcal Mf(s)=1/s\)；但 \(1/(c+\mathrm it)\) 的模只按 \(|t|^{-1}\) 衰减，不能直接使用本节的绝对反演定理。对称截断或其他正则化仍可能给出反演，需要另行陈述其收敛方式。

\Needspace{6\baselineskip}
\begin{proposition}[绝对反演的一个充分条件]\label{b3:prop:21-mellin-inversion-test}
设 \(c\in\mathbb R\)，\(f:(0,\infty)\to\mathbb C\)，且 \(g(u)=e^{cu}f(e^u)\) 属于 \(C^2(\mathbb R)\)。若 \(g,g',g''\in L^1(\mathbb R)\)，且 \(g(u),g'(u)\to0\) 当 \(u\to\pm\infty\)，则
\[
 |\mathcal Mf(c+\mathrm it)|
 \le\min\{\|g\|_1,\,|t|^{-2}\|g''\|_1\}\quad(t\ne0).
\]
因此 \(f\) 满足 Mellin 反演定理的全部条件，且反演处处成立。
\end{proposition}
\begin{proof}
范数关系给出 \(f\in L^1_c\)。在 \(\int g(u)e^{\mathrm itu}\,\mathrm du\) 中分部积分两次，端点项均因假设消失，得到模界 \(|t|^{-2}\|g''\|_1\)；不分部积分则得到 \(\|g\|_1\)。前一个界控制 \(|t|\ge1\)，后一个控制 \(|t|\le1\)，二者合并可积。\(g\) 连续也保证 \(f(x)=x^{-c}g(\log x)\) 连续，故反演处处成立。
\end{proof}

\subsection{对数坐标、乘法卷积与围道移动}
\label{b3:subsec:210505}

每条竖线对应不同的加权函数 \(g_c\)。改变 \(c\) 时，不只是给同一 Fourier 变换换了记号，而是改变了输入在两端的权重。对于伸缩结构，最自然的卷积同样应使用不变测度 \(\mathrm dt/t\)。

\Needspace{6\baselineskip}
\begin{proposition}[Mellin 卷积]\label{b3:prop:21-mellin-convolution}
设 \(c\in\mathbb R\)，可测函数 \(f,h\in L^1_c\)。则积分
\[
 (f*_M h)(x)\coloneqq\int_0^\infty f(t)h(x/t)\,\frac{\mathrm dt}{t}
\]
对几乎处处的 \(x>0\) 绝对收敛，定义 \(L^1_c\) 中的函数，并满足
\[
 \|f*_M h\|_{1,c}\le\|f\|_{1,c}\cdot\|h\|_{1,c}.
\]
在整条 \(\operatorname{Re}s=c\) 上，
\begin{equation}\label{b3:eq:21-mellin-convolution}
 \mathcal M(f*_M h)(s)=\mathcal Mf(s)\mathcal Mh(s).
\end{equation}
\end{proposition}
\begin{proof}
先对绝对值使用 Tonelli 定理，再在内部积分令 \(x=tu\)，得到
\[
 \begin{aligned}
 &\int_0^\infty\int_0^\infty
       |f(t)|\cdot|h(x/t)|\,\frac{\mathrm dt}{t}\,x^{c-1}\,\mathrm dx\\
 &=\int_0^\infty|f(t)|t^{c-1}\,\mathrm dt
        \int_0^\infty|h(u)|u^{c-1}\,\mathrm du<\infty.
 \end{aligned}
\]
这先给出卷积的几乎处处绝对存在性，再由三角不等式给出范数界。对 \(s=c+\mathrm iy\)，带复幂的相应二重积分具有同一个绝对值控制，故可以用 Fubini 交换积分。换元后 \(x^{s-1}\,\mathrm dx=t^su^{s-1}\,\mathrm du\)，外层的 \(\mathrm dt/t\) 正好将 \(t^s\) 变为 \(t^{s-1}\,\mathrm dt\)，便得到两个 Mellin 积分的乘积。
\end{proof}

这一卷积称为\term[mellin-juanji]{Mellin 卷积}[Mellin convolution]。若分别令 \(g_c(u)=e^{cu}f(e^u)\)、\(k_c(u)=e^{cu}h(e^u)\)，则直接换元给
\[
 e^{cu}(f*_Mh)(e^u)=(g_c*k_c)(u)
\]
几乎处处成立。乘法卷积就是对数坐标下的加法卷积；若把 \(\mathrm dt/t\) 错换成 \(\mathrm dt\)，权重和乘积公式都会随之改变。

\ProseParagraphBreak

移动 Gamma 反演的围道时，水平连接边扫过有限实部条带，因此需要在整个条带上统一控制高处的函数值。上一节特殊竖线的精确模值提示了衰减现象；从 Euler 积分出发的分部积分将给出移线所需的统一衰减。

\Needspace{6\baselineskip}
\begin{lemma}[Gamma 函数在有限条带的竖向衰减]\label{b3:lem:21-gamma-strip-decay}
设 \(a<b\) 为实数，\(k\ge0\) 为整数。存在常数 \(C_{a,b,k}\)，使 Gamma 函数的亚纯延拓满足
\[
 |\Gamma(\sigma+\mathrm it)|\le C_{a,b,k}(1+|t|)^{-k},
 \quad a\le\sigma\le b,\quad |t|\ge1.
\]
\end{lemma}
\begin{proof}
先设 \(a>0\)。Euler 积分令 \(x=e^u\)，得到
\[
 \Gamma(\sigma+\mathrm it)
 =\int_{\mathbb R}g_\sigma(u)e^{\mathrm itu}\,\mathrm du,
 \quad g_\sigma(u)=e^{\sigma u-e^u}.
\]
对每个固定 \(j\)，逐次求导可写成
\[
 g_\sigma^{(j)}(u)=g_\sigma(u)P_j(\sigma,e^u),
\]
其中 \(P_j\) 是多项式。事实上 \(P_0=1\)，递推为
\[
 P_{j+1}(\sigma,v)=(\sigma-v)P_j(\sigma,v)
                         +v\frac{\partial P_j}{\partial v}(\sigma,v).
\]
\(\sigma\in[a,b]\) 时系数统一有界。在 \(u\le0\) 上，\(e^u\le1\)，所以每个导数的模由 \(Ce^{au}\) 控制；在 \(u\ge0\) 上，可由 \(Ce^{(b+j)u-e^u}\) 控制。两端均可积，且各阶导数的边界值都为零。分部积分 \(k\) 次，便得到统一的 \(C|t|^{-k}\) 上界。

一般 \(a\) 下，选正整数 \(q\) 使 \(a+q>0\)。递推给
\[
 \Gamma(\sigma+\mathrm it)
 =\frac{\Gamma(\sigma+q+\mathrm it)}
         {\prod_{j=0}^{q-1}(\sigma+j+\mathrm it)}.
\]
当 \(|t|\ge1\) 时，每个分母因子的模至少为 \(|t|\)，分子已有正实部闭带上的任意固定幂次衰减。故所需估计也对整个给定条带成立。所有极点都在实轴上，当前的高处区域不遇到它们。
\end{proof}

\(e^{-x}\in L^1_c\) 对每个 \(c>0\) 成立，而引理使 \(t\mapsto\Gamma(c+\mathrm it)\) 绝对可积。因此由反演定理，
\begin{equation}\label{b3:eq:21-exponential-inversion}
 e^{-x}=\frac1{2\pi\mathrm i}
       \int_{c-\mathrm i\infty}^{c+\mathrm i\infty}\Gamma(s)x^{-s}\,\mathrm ds,
 \quad c>0,\quad x>0.
\end{equation}
这一步已经独立恢复原函数。下面的留数计算将进一步恢复它的 Taylor 级数，不能把两个步骤互相作为未经证明的依据。

\Needspace{6\baselineskip}
\begin{proposition}[指数函数反演的移线余项]\label{b3:prop:21-exponential-contour}
对每个整数 \(N\ge0\) 和实数 \(x>0\)，有
\[
 e^{-x}=\sum_{n=0}^N\frac{(-x)^n}{n!}+R_N(x),
\]
其中
\begin{equation}\label{b3:eq:21-exponential-remainder}
 R_N(x)=\frac1{2\pi\mathrm i}
       \int_{-N-1/2-\mathrm i\infty}^{-N-1/2+\mathrm i\infty}
                    \Gamma(s)x^{-s}\,\mathrm ds
\end{equation}
绝对收敛。存在与 \(N,x\) 无关的常数 \(C\)，使
\[
 |R_N(x)|\le C\frac{x^{N+1/2}}{\Gamma(N+3/2)}.
\]
因此余项在每个固定区间 \(0<x\le X<\infty\) 上一致趋于零。
\end{proposition}
\begin{proof}
在\eqref{b3:eq:21-exponential-inversion}中取 \(c=1\)。固定 \(N,x\)，考虑左右边实部分别为 \(-N-1/2\) 与 \(1\)、上下边虚部分别为 \(\pm T\) 的矩形，按正向绕行。其右边由下到上，左边由上到下。矩形内的极点恰为 \(0,-1,\ldots,-N\)，由\secref{b3:sec:2102}的留数计算，
\[
 \operatorname{Res}_{s=-n}\bigl(\Gamma(s)x^{-s}\bigr)
 =\frac{(-1)^nx^n}{n!}.
\]
对水平边使用\cref{b3:lem:21-gamma-strip-decay}，取 \(k=2\)。在这个固定实部区间，\(|x^{-s}|=x^{-\operatorname{Re}s}\) 有界，两条水平边的总长固定，所以其积分为 \(O_{N,x}(T^{-2})\)，趋于零。两个竖线积分都绝对收敛，故令 \(T\to\infty\)，留数定理给出所述有限和与余项，符号由上述边界方向确定。

为同时控制 \(N\)，将新竖线上的 Gamma 向右递推 \(N+1\) 次：
\[
 |\Gamma(-N-1/2+\mathrm it)|
 \le\frac{|\Gamma(1/2+\mathrm it)|}
              {\prod_{j=0}^N(j+1/2)}
 =\frac{\sqrt\pi}{\Gamma(N+3/2)}
                         |\Gamma(1/2+\mathrm it)|.
\]
上一节的精确模值 \(|\Gamma(1/2+\mathrm it)|^2=\pi/\cosh\pi t\) 表明右端的分子关于 \(t\) 可积，积分为一个固定常数。又新直线上 \(|x^{-s}|=x^{N+1/2}\)，所以得到所示统一余项界。

固定 \(X>0\)，令 \(a_N=X^{N+1/2}/\Gamma(N+3/2)\)。递推给
\[
 \frac{a_{N+1}}{a_N}=\frac{X}{N+3/2}\longrightarrow0,
\]
故 \(a_N\to0\)。对所有 \(0<x\le X\) 以 \(a_N\) 控制即可得到一致性。这一步甚至不需要再次使用 Stirling 公式。
\end{proof}

最后把移线估计中真正使用的条件抽出来。它们与函数是否叫作 Gamma 无关，也说明为什么仅有亚纯延拓还不够。

\Needspace{6\baselineskip}
\begin{proposition}[有限条带内的 Mellin 移线]\label{b3:prop:21-mellin-contour-shift}
设 \(a<b\) 为实数，\(F\) 在闭条带 \(a\le\operatorname{Re}s\le b\) 的某个开邻域亚纯，在该闭条带内的极点恰为有限个点 \(s_1,\ldots,s_m\)，且边界两条竖线上无极点。假设
\[
 \int_{\mathbb R}|F(a+\mathrm it)|\,\mathrm dt<\infty,
 \quad\int_{\mathbb R}|F(b+\mathrm it)|\,\mathrm dt<\infty,
\]
并存在 \(T_j\to\infty\)，使
\[
 \sup_{a\le\sigma\le b}|F(\sigma+\mathrm iT_j)|
 +\sup_{a\le\sigma\le b}|F(\sigma-\mathrm iT_j)|\longrightarrow0.
\]
则对每个 \(x>0\)，两个竖线积分绝对收敛，并满足
\begin{equation}\label{b3:eq:21-mellin-contour-shift}
 \frac1{2\pi\mathrm i}\int_{b-\mathrm i\infty}^{b+\mathrm i\infty}F(s)x^{-s}\,\mathrm ds
 =\sum_{\nu=1}^m\operatorname{Res}_{s=s_\nu}\bigl(F(s)x^{-s}\bigr)
  +\frac1{2\pi\mathrm i}\int_{a-\mathrm i\infty}^{a+\mathrm i\infty}F(s)x^{-s}\,\mathrm ds.
\end{equation}
若 \(F\) 在 \(s_0\) 的主要部分为 \(\sum_{k=1}^q A_k(s-s_0)^{-k}\)，则相应留数为
\[
 \operatorname{Res}_{s=s_0}\bigl(F(s)x^{-s}\bigr)
 =x^{-s_0}\sum_{k=1}^q\frac{A_k(-\log x)^{k-1}}{(k-1)!}.
\]
\end{proposition}
\begin{proof}
取上述两条竖线和高度 \(\pm T_j\) 的正向矩形。\(j\) 充分大时，它包含全部指定极点。水平边上 \(|x^{-s}|\) 至多为 \(\max\{x^{-a},x^{-b}\}\)，所以每条水平边的积分模至多为
\[
 (b-a)\max\{x^{-a},x^{-b}\}
           \sup_{a\le\sigma\le b}|F(\sigma\pm\mathrm iT_j)|,
\]
其极限为零。竖线上的绝对可积性使截断积分趋于整个直线积分。留数定理于是给出\eqref{b3:eq:21-mellin-contour-shift}。

在 \(s_0\) 附近展开
\[
 x^{-s}=x^{-s_0}\sum_{j=0}^{\infty}
                   \frac{(-\log x)^j(s-s_0)^j}{j!}.
\]
与有限主要部分相乘，选取 \((s-s_0)^{-1}\) 的系数，就得到最后的公式。
\end{proof}

例如 \(s=-\alpha\) 处的 \(q\) 阶极点贡献 \(x^\alpha\) 乘一个次数至多 \(q-1\) 的 \(\log x\) 多项式，与\cref{b3:prop:21-mellin-endpoint}中的正向计算相互核对。若把直线向左移到 \(a\)，新积分的模至多为
\[
 \frac{x^{-a}}{2\pi}\int_{\mathbb R}|F(a+\mathrm it)|\,\mathrm dt.
\]
这给 \(x\downarrow0\) 时可读出的余项尺度；向右移线则控制 \(x\to\infty\)，并应按\eqref{b3:eq:21-mellin-contour-shift}重新移项，保留留数符号。端点展开、亚纯极点和反演渐近由此组成一条有条件的双向联系，后续 theta 与 zeta 的积分方法会再次使用它。


\begin{chaptersummary}
Euler 积分先在右半平面定义 Gamma 函数，其全纯性、参数导数与递推都来自可以逐项核对的两端控制。正实轴上的严格对数凸性不仅排除非平凡周期因子，还给出固定平移的比值估计。Bohr--Mollerup 的条件共同确定函数，单独的阶乘插值或递推并不足够。

\ProseParagraphBreak
递推延拓与 Taylor 扣除描述同一个亚纯函数。全部非正整数都是单极点，留数决定奇部，有限部分进一步决定与零点相消后的导数。倒数乘积则说明 Gamma 没有复零点；其线性指数常数、对数导数与局部幂和都由 Euler 极限统一确定。

\ProseParagraphBreak
反射、Beta、倍增和乘法公式分别利用相反参数、积分换元与递推唯一性，规范化常数必须随这些步骤一起保留。Stirling 公式来自大参数下积分质量的集中；固定平移、增长平移与竖向参数是不同的极限问题，不能把一个区域中的渐近无条件移用到另一个区域。

\ProseParagraphBreak
Mellin 变换在对数坐标中成为带权 Fourier 变换，缩放成为乘子，乘法卷积成为普通乘积。收敛带记录原函数两端的可积性，亚纯延拓提供更多参数值，却不改变原积分的收敛范围。绝对反演、水平边消失和远处竖线余项趋零分别支持恢复函数、有限移线和最终留数展开；端点幂次决定极点位置，对数幂决定极点阶数。
\end{chaptersummary}

\BookExercises

% 自拟：由Gamma对数矩进入最优仿射逼近，利用协方差和正核表示求出严格剩余误差及其精确尺度。
\Needspace{6\baselineskip}
\begin{exercise}[对数函数的最优仿射逼近]\label{b3:ex:21-1}
对 \(x>0\)，用
\[
 \mathbb E_x[H(T)]=\frac1{\Gamma(x)}
 \int_0^\infty H(t)t^{x-1}e^{-t}\,\mathrm{d}t
\]
表示该正权重下的平均，并令 \(\psi_0=\Gamma'/\Gamma\)。在所有实常数 \(A,B\) 中，求使
\[
 \mathbb E_x\bigl[(\log T-AT-B)^2\bigr]
\]
最小的唯一一对 \((A,B)\)。证明最小值为
\(D(x)=\psi_0'(x)-1/x>0\)，并进一步证明
\[
 \lim_{x\to\infty}x^2D(x)=\frac12.
\]
这里允许使用本章的对数导数级数，但须说明最小值为何严格为正，以及最后一个极限的支配条件。
\end{exercise}
\begin{solution}
Euler 积分的递推与参数微分给
\[
 \mathbb E_x[T]=x,\quad
 \mathbb E_x[T^2]=x(x+1),\quad
 \mathbb E_x[\log T]=\psi_0(x),\quad
 \operatorname{Var}_x(\log T)=\psi_0'(x).
\]
这里方差表示 \(\mathbb E_x[(Y-\mathbb E_xY)^2]\)。所需混合矩同样绝对可积；原点附近的幂乘对数可积，无穷远仍有指数下降。于是
\[
 \mathbb E_x[T\log T]
 =\frac{\Gamma'(x+1)}{\Gamma(x)}
 =x\psi_0(x)+1,
\]
最后一步由 \(\Gamma'(x+1)=\Gamma(x)+x\Gamma'(x)\) 得到。因此 \(T\) 的方差为 \(x\)，而 \(T\) 与 \(\log T\) 的协方差为一。

把两变量中心化并配方，得到精确恒等式
\[
 \mathbb E_x\bigl[(\log T-AT-B)^2\bigr]
 =\psi_0'(x)-\frac1x
   +x\left(A-\frac1x\right)^2
   +\bigl(B+Ax-\psi_0(x)\bigr)^2.
\]
故唯一最优选择为
\(A=1/x\)、\(B=\psi_0(x)-1\)。若相应最小值为零，连续函数 \(\log t-t/x-B\) 必须在这个处处正的密度下几乎处处为零，继而在全部 \(t>0\) 为零。但其导数 \(1/t-1/x\) 不恒为零，矛盾，所以 \(D(x)>0\)。

为了求精确极限，由正项级数和 Tonelli 定理有
\[
 \psi_0'(x)=\sum_{n=0}^\infty\frac1{(x+n)^2}
 =\int_0^\infty\frac{u e^{-xu}}{1-e^{-u}}\,\mathrm{d}u.
\]
减去 \(1/x=\int_0^\infty e^{-xu}\,\mathrm{d}u\)，记
\(q(u)=u/(1-e^{-u})-1\)。对 \(u>0\)，不等式
\(1-e^{-u}<u\) 和 \(e^u>1+u\) 分别给 \(0<q(u)<u\)；Taylor 展开还给 \(q(u)=u/2+O(u^2)\)。令 \(v=xu\)，则
\[
 x^2D(x)=\int_0^\infty e^{-v}xq(v/x)\,\mathrm{d}v.
\]
被积函数逐点趋于 \(ve^{-v}/2\)，并由可积函数 \(ve^{-v}\) 控制。支配收敛遂给极限 \(\Gamma(2)/2=1/2\)。这个剩余误差衡量的是去掉最佳线性部分后的对数波动。
\end{solution}

% 自拟：把阻尼Gamma积分的计算与统一振荡尾估计相连，区分去阻尼极限和通常广义积分，并验证条件而非绝对收敛。
\Needspace{6\baselineskip}
\begin{exercise}[去阻尼与广义积分]\label{b3:ex:21-2}
设 \(b,c>0\)、\(\varepsilon>0\)，计算
\[
 I_\varepsilon(c,b)=\int_0^\infty
 x^{c-1}e^{-\varepsilon x}\cos(bx)\,\mathrm{d}x,
\]
并求 \(\varepsilon\downarrow0\) 的极限。证明当 \(0<c<1\) 时，通常的广义积分
\(\int_0^\infty x^{c-1}\cos(bx)\,\mathrm{d}x\) 条件收敛，且等于这一极限。再以 \(c=1\) 说明：去阻尼极限存在并不保证通常广义积分存在。
\end{exercise}
\begin{solution}
对 \(\operatorname{Re}w>0\)，函数
\[
 J(w)=\int_0^\infty x^{c-1}e^{-wx}\,\mathrm{d}x
\]
全纯；在每个紧参数集上，积分由 \(x^{c-1}e^{-\delta x}\) 控制，参数导数加入的幂也可积。正实 \(w\) 上换元给 \(J(w)=\Gamma(c)w^{-c}\)。在右半平面取正实轴上为实值的对数分支，恒等定理把该式推广到整个右半平面。令 \(w=\varepsilon-\mathrm{i}b\)，取实部，得
\[
 I_\varepsilon(c,b)
 =\Gamma(c)(\varepsilon^2+b^2)^{-c/2}
 \cos\left(c\arctan\frac b\varepsilon\right).
\]
所以对每个 \(c>0\)，去阻尼极限都为
\(\Gamma(c)b^{-c}\cos(\pi c/2)\)。

现在设 \(0<c<1\)。对 \(\varepsilon\ge0\)，振幅
\(a_\varepsilon(x)=x^{c-1}e^{-\varepsilon x}\) 在 \(x>0\) 单调递减并趋于零。分部积分给共同尾界
\[
 \left|\int_R^\infty a_\varepsilon(x)e^{\mathrm{i}bx}\,\mathrm{d}x\right|
 \le\frac{a_\varepsilon(R)+\int_R^\infty|a_\varepsilon'(x)|\,\mathrm{d}x}{b}
 =\frac{2a_\varepsilon(R)}b
 \le\frac{2R^{c-1}}b.
\]
当 \(\varepsilon=0\) 时，右侧同时证明振荡尾积分存在。在固定区间 \((0,R)\)，以 \(x^{c-1}\) 支配即可令 \(\varepsilon\downarrow0\)；再令 \(R\to\infty\)，上述统一尾界使两个极限相容。因此
\[
 \int_0^\infty x^{c-1}\cos(bx)\,\mathrm{d}x
 =\Gamma(c)b^{-c}\cos\frac{\pi c}2.
\]
这不是绝对收敛。围绕 \(2\pi k/b\) 取固定长度的小区间，使其中 \(|\cos(bx)|\ge1/2\)。第 \(k\) 个区间上的绝对积分不小于 \(Ck^{c-1}\)，而该正项级数发散。

最后取 \(c=1\)。已得公式给 \(I_\varepsilon(1,b)=\varepsilon/(\varepsilon^2+b^2)\to0\)，但
\(\int_0^R\cos(bx)\,\mathrm{d}x=\sin(bR)/b\) 没有极限。前面的统一尾界正是在 \(c=1\) 时失去趋零性质，因而不能把同一论证用到端点。
\end{solution}

% 自拟：将任意平移Gamma商的零极点分类与共同极点处的有限部分相连，辨明何时真正退化为有理函数。
\Needspace{6\baselineskip}
\begin{exercise}[Gamma 商的零极点与相消]\label{b3:ex:21-3}
设 \(a,b\in\mathbb C\)，把
\[
 Q(z)=\frac{\Gamma(z+a)}{\Gamma(z+b)}
\]
理解为亚纯函数。确定它的全部零点和极点，并分别判定它何时是有理函数、多项式或某个多项式的倒数。若某点 \(z_0\) 满足
\(z_0+a=-n\)、\(z_0+b=-r\)，其中 \(n,r\) 为非负整数，求 \(Q\) 在这个共同极点处可去延拓后的值与导数。
\end{exercise}
\begin{solution}
Gamma 没有零点，而其倒数在非正整数有且仅有单零点。因此若 \(b-a\notin\mathbb Z\)，两列点不相交，\(Q\) 的极点恰为 \(-a-n\)，零点恰为 \(-b-n\)，其中 \(n\ge0\)，且均为一阶。没有其他零极点。非零有理函数只有有限多个零极点，所以此时 \(Q\) 不是有理函数。

若 \(b-a=m\) 是正整数，有限次递推给
\[
 Q(z)=\frac1{\prod_{k=0}^{m-1}(z+a+k)}.
\]
故只剩 \(-a,-a-1,\ldots,-a-m+1\) 这些单极点，没有零点。若 \(a-b=m\) 是正整数，则
\[
 Q(z)=\prod_{k=0}^{m-1}(z+b+k),
\]
其零点为 \(-b,-b-1,\ldots,-b-m+1\)，均为单零点，没有极点。若 \(a=b\)，则 \(Q=1\)。这些恒等式先在避开极点的开集成立，再作为亚纯恒等式延拓，已经包含所有表观奇点的相消。

所以 \(Q\) 为有理函数当且仅当 \(a-b\in\mathbb Z\)；为多项式当且仅当 \(a-b\) 是非负整数；为非零多项式的倒数当且仅当 \(b-a\) 是非负整数。常函数一同时属于后两类。

最后令 \(w=z-z_0\)。记 \(H_0=0\)、\(H_j=\sum_{k=1}^j1/k\)。本章的极点有限部分给
\[
 \Gamma(-n+w)=\frac{(-1)^n}{n!w}
 \left(1+(H_n-\gamma)w+O(w^2)\right),
\]
分母有把 \(n\) 换为 \(r\) 的同类展开。两式相除得
\[
 Q(z_0+w)=(-1)^{n-r}\frac{r!}{n!}
 \left(1+(H_n-H_r)w+O(w^2)\right).
\]
因此
\[
 Q(z_0)=(-1)^{n-r}\frac{r!}{n!},\quad
 Q'(z_0)=(-1)^{n-r}\frac{r!}{n!}(H_n-H_r).
\]
值由留数之比确定，导数还需要两边有限部分的差；形式上的“两个无穷大相除”不能给出这些资料。
\end{solution}

% 自拟：构造保持递推、规范化与严格普通凸性的非平凡周期扰动，说明Bohr--Mollerup中的对数凸性不可直接弱化。
\Needspace{6\baselineskip}
\begin{exercise}[普通凸性不能代替对数凸性]\label{b3:ex:21-4}
对实数 \(\varepsilon\)，令
\[
 F_\varepsilon(x)=\Gamma(x)e^{\varepsilon\sin2\pi x},\quad x>0.
\]
证明存在 \(\varepsilon_0>0\)，使每个 \(0<|\varepsilon|\le\varepsilon_0\) 对应的 \(F_\varepsilon\) 都正且严格凸，满足
\(F_\varepsilon(1)=1\) 与 \(F_\varepsilon(x+1)=xF_\varepsilon(x)\)，却不等于 \(\Gamma\)。再证明只要 \(\varepsilon\ne0\)，\(\log F_\varepsilon\) 就不是凸函数。无需给出 \(\varepsilon_0\) 的数值，但须证明一个对所有 \(x>0\) 有效的选择确实存在。
\end{exercise}
\begin{solution}
正性、规范化与递推直接来自指数因子的一周期性。若 \(\varepsilon\ne0\)，取 \(x=1/4\) 即知这个因子不等于一，所以 \(F_\varepsilon\ne\Gamma\)。问题在于选择一个足够小而对全半轴适用的扰动。

置 \(\omega=2\pi\)，并在本解中记 \(\psi=\psi_0\)。对数导数级数给 \(\psi'(x)>0\)；原点展开给 \(\psi(x)=-1/x+O(1)\)，而
\(\psi(n)=H_{n-1}-\gamma\) 与单调性给 \(\psi(x)\to\infty\)。因此 \(|\psi(x)|\) 在 \(x\downarrow0\) 和 \(x\to\infty\) 两端都趋于无穷。

定义连续正函数
\[
 q(x)=\psi'(x)+\psi(x)^2,\quad
 r(x)=2\omega|\psi(x)|+\omega^2.
\]
有
\[
 \frac{q(x)}{r(x)}
 \ge\frac{\psi(x)^2}{2\omega|\psi(x)|+\omega^2}
 \longrightarrow\infty
\]
于两端。在中间紧区间上，这个连续正商取得正的最小值，故
\(\mu=\inf_{x>0}q(x)/r(x)>0\)。取 \(\varepsilon_0=\mu/2\)。逐次求导得
\[
 \frac{F_\varepsilon''}{F_\varepsilon}
 =q+2\varepsilon\omega\psi\cos(\omega x)
   -\varepsilon\omega^2\sin(\omega x)
   +\varepsilon^2\omega^2\cos^2(\omega x)
 \ge q-|\varepsilon|r\ge\frac q2>0.
\]
因此所有这样的 \(F_\varepsilon\) 都严格凸。

另一方面，
\[
 (\log F_\varepsilon)''(x)
 =\psi'(x)-\varepsilon\omega^2\sin(\omega x).
\]
本章的级数积分比较给 \(\psi'(x)\to0\)。当 \(\varepsilon>0\) 时取 \(x_n=n+1/4\)，当 \(\varepsilon<0\) 时取 \(x_n=n+3/4\)，两种情形都有
\[
 (\log F_\varepsilon)''(x_n)
 =\psi'(x_n)-|\varepsilon|\omega^2<0
\]
于充分大的 \(n\)。所以对数不凸。普通凸性可以由一阶对数导数的平方维持，而对数凸性只看二阶对数导数；这正是本例中两个条件出现不同结论的原因。
\end{solution}

% 自拟：给同一有理亚纯函数选取三个不同反演竖线，求出各自原函数与实际收敛带，显示移线留数的作用。
\Needspace{6\baselineskip}
\begin{exercise}[同一表达式的三个反演条带]\label{b3:ex:21-5}
设 \(\alpha<\beta\) 为实数，令
\[
 F(s)=\frac1{(s-\alpha)(\beta-s)}.
\]
对 \(x>0\) 及 \(c\in\mathbb R\setminus\{\alpha,\beta\}\)，证明
\[
 I_c(x)=\frac1{2\pi\mathrm{i}}
 \int_{c-\mathrm{i}\infty}^{c+\mathrm{i}\infty}F(s)x^{-s}\,\mathrm{d}s
\]
绝对收敛，并分别计算 \(c<\alpha\)、\(\alpha<c<\beta\) 和 \(c>\beta\) 时的函数。求这三个函数各自的最大绝对收敛带，用留数解释相邻两种答案的差，并说明为什么它们不违背 Mellin 变换的唯一性。
\end{exercise}
\begin{solution}
在任意不经过极点的竖线上，\(F(c+\mathrm{i}t)=O((1+|t|)^{-2})\)，而 \(|x^{-c-\mathrm{i}t}|=x^{-c}\)，故所示积分绝对收敛。记 \(d=\beta-\alpha>0\)，考虑三个连续函数
\[
 f_0(x)=\frac1d
 \begin{cases}
 x^{-\alpha},&0<x\le1,\\
 x^{-\beta},&x>1,
 \end{cases}
\]
以及
\[
 f_-(x)=\frac1d
 \begin{cases}
 0,&0<x\le1,\\
 x^{-\beta}-x^{-\alpha},&x>1,
 \end{cases}
 \quad
 f_+(x)=\frac1d
 \begin{cases}
 x^{-\alpha}-x^{-\beta},&0<x<1,\\
 0,&x\ge1.
 \end{cases}
\]
对 \(f_0\)，零端要求 \(\operatorname{Re}s>\alpha\)，无穷远要求 \(\operatorname{Re}s<\beta\)，且
\[
 \mathcal Mf_0(s)=\frac1d
 \left(\frac1{s-\alpha}+\frac1{\beta-s}\right)=F(s).
\]
对 \(f_-\)，只有无穷远端需要检查，那里
\(f_-(x)\sim-x^{-\alpha}/d\)。其最大绝对收敛域为 \(\operatorname{Re}s<\alpha\)，在该域逐项积分给
\[
 \mathcal Mf_-(s)=\frac1d
 \left(\frac1{\beta-s}-\frac1{\alpha-s}\right)=F(s).
\]
同样，\(f_+(x)\sim-x^{-\beta}/d\) 于零端，其最大绝对收敛域为 \(\operatorname{Re}s>\beta\)，并且
\[
 \mathcal Mf_+(s)=\frac1d
 \left(\frac1{s-\alpha}-\frac1{s-\beta}\right)=F(s).
\]
在各边界实部，绝对积分都含非零常数倍的 \(\int\mathrm{d}x/x\)，所以不能把边界加入。三个函数都在各自竖线的权重下可积，变换沿竖线也可积；Mellin 绝对反演因此处处给
\[
 I_c=
 \begin{cases}
 f_-,&c<\alpha,\\
 f_0,&\alpha<c<\beta,\\
 f_+,&c>\beta.
 \end{cases}
\]
在 \(x=1\) 时，中间答案为 \(1/d\)，两侧答案为零，这些值也由各函数的连续性确定。

反演被积函数在两个极点的留数为
\[
 \operatorname{Res}_{s=\alpha}\bigl(F(s)x^{-s}\bigr)=\frac{x^{-\alpha}}d,
 \quad
 \operatorname{Res}_{s=\beta}\bigl(F(s)x^{-s}\bigr)=-\frac{x^{-\beta}}d.
\]
由下往上的竖线向右移时，右线积分等于左线积分加所越过极点的留数。因此
\[
 f_0-f_-=\frac{x^{-\alpha}}d,\quad
 f_+-f_0=-\frac{x^{-\beta}}d.
\]
有限水平区间上 \(F(\sigma+\mathrm{i}T)=O(T^{-2})\)，且 \(x^{-\sigma}\) 一致有界，所以矩形的水平边确实趋于零。唯一性要求两个原函数在同一权重下可积，并在共同竖线上有相同变换；这里三个最大收敛带互不相交，不满足这个前提。只保留亚纯表达式而忘记收敛带，会丢失原函数信息。
\end{solution}

% 自拟：以高阶有限差分制造原点抵消，计算被消去的Gamma极点处的值，并用正积分证明对应零点确为单零点。
\Needspace{6\baselineskip}
\begin{exercise}[高阶抵消与可去值]\label{b3:ex:21-6}
设 \(a>0\)，\(m\ge1\) 为整数，令
\[
 f_m(x)=e^{-ax}(1-e^{-x})^m,\quad
 A_m(s)=\sum_{k=0}^m(-1)^k\binom mk(a+k)^{-s},
\]
其中幂使用正实数的实对数。求 \(f_m\) 的最大 Mellin 绝对收敛域，并证明其中
\[
 \mathcal Mf_m(s)=\Gamma(s)A_m(s)
\]
按可去延拓解释。计算 \(s=-j\)、\(0\le j<m\) 处的值，证明这些点都是 \(A_m\) 的恰好一重零点；最后求右端在 \(s=-m\) 的留数。
\end{exercise}
\begin{solution}
在原点附近，\(f_m(x)=x^m(1+O(x))\)；无穷远则由 \(e^{-ax}\) 控制。因此最大绝对收敛域为 \(\operatorname{Re}s>-m\)，边界实部已有 \(x^{-1}\) 型发散。在这个半平面上，Mellin 积分定义全纯函数，记为 \(G_m\)。

在较小的半平面 \(\operatorname{Re}s>0\)，有限二项展开后的每项都可以单独积分，所以
\[
 G_m(s)=\sum_{k=0}^m(-1)^k\binom mk
 \int_0^\infty e^{-(a+k)x}x^{s-1}\,\mathrm{d}x
 =\Gamma(s)A_m(s).
\]
\(A_m\) 是整函数。要核对前 \(m\) 个 Gamma 极点的抵消，使用有限恒等式
\[
 \sum_{k=0}^m(-1)^k\binom mk e^{(a+k)t}
 =e^{at}(1-e^t)^m=(-1)^mt^m+O(t^{m+1}).
\]
在原点逐次求导便得
\[
 \sum_{k=0}^m(-1)^k\binom mk(a+k)^j=0\quad(0\le j<m),
\]
而 \(j=m\) 时该和等于 \((-1)^m m!\)。因此
\(A_m(-j)=0\) 对 \(0\le j<m\) 成立，足以消去这些点的单极点。恒等定理遂把 \(G_m=\Gamma A_m\) 推广到整个 \(\operatorname{Re}s>-m\)。

在 \(s=-j\) 附近，Gamma 的留数为 \((-1)^j/j!\)，故可去值为
\[
 G_m(-j)=\frac{(-1)^j}{j!}A_m'(-j)
 =\frac{(-1)^{j+1}}{j!}
 \sum_{k=0}^m(-1)^k\binom mk(a+k)^j\log(a+k).
\]
这里 \(-j>-m\)，原来的正积分仍收敛，而且 \(f_m(x)>0\)，所以 \(G_m(-j)>0\)。于是 \(A_m'(-j)\ne0\)，证明这些零点恰为一重。零点阶数不能仅从前面的有限差分恒等式直接读出，正积分排除了更高阶相消。

在 \(s=-m\)，已经算得 \(A_m(-m)=(-1)^m m!\ne0\)，因此右端仍有单极点，留数为
\[
 \frac{(-1)^m}{m!}A_m(-m)=1.
\]
它与原点展开的首项系数一相符。有限差分消除了前 \(m\) 个极点，也恰好把原积分的左收敛边界推进了 \(m\) 个单位。
\end{solution}

% 自拟：从不同指数截断幂的乘法卷积导出部分分式公式，再通过紧单纯形积分解释参数合并时出现的对数因子。
\Needspace{6\baselineskip}
\begin{exercise}[不同指数的卷积与合并极限]\label{b3:ex:21-7}
设 \(m\ge2\)，\(a_1,\ldots,a_m\) 为实数，并定义
\[
 h_a(x)=
 \begin{cases}
 x^a,&0<x<1,\\
 0,&x\ge1,
 \end{cases}
 \quad H=h_{a_1}*_M\cdots*_Mh_{a_m},
\]
其中 \((f*_Mg)(x)=\int_0^\infty f(t)g(x/t)\,\mathrm{d}t/t\)。证明当参数两两不同时，
\[
 H(x)=
 \begin{cases}
 \displaystyle\sum_{j=1}^m
 \frac{x^{a_j}}{\prod_{k\ne j}(a_k-a_j)},&0<x<1,\\
 0,&x\ge1.
 \end{cases}
\]
再证明当所有 \(a_j\to a\) 时，这个表达式在 \((0,\infty)\) 的紧集上一致趋于
\[
 \frac{x^a(-\log x)^{m-1}}{(m-1)!}\cdot\mathbf1_{(0,1)}(x).
\]
求两种情形的 Mellin 变换与最大绝对收敛域，并说明系数分母趋零为什么不导致卷积发散。
\end{exercise}
\begin{solution}
先用对数坐标写出不含相消的积分。令 \(x=e^{-u}\)。当 \(u>0\) 时，多次卷积和换元给
\[
 H(e^{-u})=u^{m-1}\int_{\Delta}
 \exp\left[-u\left(\sum_{j=1}^{m-1}a_jw_j
              +a_m\left(1-\sum_{j=1}^{m-1}w_j\right)\right)\right]
 \,\mathrm{d}w_1\cdots\mathrm{d}w_{m-1},
\]
其中
\[
 \Delta=\left\{\,(w_1,\ldots,w_{m-1})\in[0,\infty)^{m-1}
          \ \middle|\ \sum_{j=1}^{m-1}w_j\le1\,\right\}.
\]
具体地，两因子卷积在此坐标变成 \(\int_0^u\) 上的加法卷积；归纳后令各积分变量为 \(uw_j\)，就得到上式。所有被积函数非负，Tonelli 定理保证迭代次序合法。若 \(x\ge1\)，支集条件没有正长度的积分区域，故 \(H(x)=0\)。这个紧单纯形表示还证明 \(H\) 连续，且因 \(m\ge2\)，在 \(x=1\) 的值为零。

设 \(a_* =\min_j a_j\)。在 \(\operatorname{Re}s>-a_*\)，各因子都在同一权重下可积，卷积定理给
\[
 \mathcal MH(s)=\prod_{j=1}^m\frac1{s+a_j}.
\]
当参数两两不同时，部分分式展开为
\[
 \prod_{j=1}^m\frac1{s+a_j}
 =\sum_{j=1}^m\frac{c_j}{s+a_j},\quad
 c_j=\frac1{\prod_{k\ne j}(a_k-a_j)}.
\]
系数由乘以 \(s+a_j\) 后令 \(s=-a_j\) 得到。每项正是 \(c_jh_{a_j}\) 的 Mellin 变换，因此题给候选式与 \(H\) 在共同竖线上有相同变换。该有理函数沿竖线为 \(O(|t|^{-m})\)，属于 \(L^1\)。此外 \(\sum_jc_j=0\)，因为在部分分式两侧令 \(s\to+\infty\) 后比较 \(1/s\) 项即可；故候选式在 \(x=1\) 也连续。对两个连续函数使用绝对反演，便知它们处处相等。

若最小参数为 \(a_*=a_{j_*}\)，其系数
\(c_{j_*}=\prod_{k\ne j_*}(a_k-a_*)^{-1}>0\)，其余幂次都较大。因此
\(H(x)\sim c_{j_*}x^{a_*}\) 于零端，最大绝对收敛域确为 \(\operatorname{Re}s>-a_*\)；无穷远没有贡献。

现在让所有 \(a_j\to a\)。紧单纯形上的被积函数在每个有界 \(u\ge0\) 区间上一致趋于 \(e^{-au}\)。单纯形的体积为 \(1/(m-1)!\)，可由逐次积分归纳得到，因而
\[
 H(e^{-u})\longrightarrow
 \frac{e^{-au}u^{m-1}}{(m-1)!}
\]
在这些区间上一致成立。任意 \(x>0\) 的紧区间都只对应有界的 \(u\) 区间；\(x\ge1\) 时又恒为零，所以得到题中的局部一致极限。在参数全部等于 \(a\) 时，上式本身就是精确公式，且
\[
 \mathcal MH(s)=\frac1{(s+a)^m},\quad
 \operatorname{Re}s>-a.
\]
零端的 \(x^a(-\log x)^{m-1}\) 说明这个收敛域同样最大。分母 \(\prod_{k\ne j}(a_k-a_j)\) 的趋零只是部分分式表达中的相消变得显著；原卷积的紧区域积分始终有界并连续依赖参数。简单极点合并成高阶极点，原函数中相应出现对数幂。
\end{solution}

% 自拟：把固定平移渐近推广到增长参数的一致区间，计算二次指数修正的有效尺度，并在临界尺度保留三次项。
\Needspace{6\baselineskip}
\begin{exercise}[增长平移的有效尺度]\label{b3:ex:21-8}
对实数 \(x>0\)、\(h\) 满足 \(x+h>0\)，令
\[
 R(x,h)=\frac{\Gamma(x+h)}{x^h\Gamma(x)}.
\]
设 \(L(x)\ge0\) 且 \(L(x)=o(x^{2/3})\)。证明
\[
 \sup_{|h|\le L(x)}
 \left|\log R(x,h)-\frac{h^2}{2x}\right|
 \longrightarrow0\quad(x\to\infty).
\]
由此分别讨论 \(h=o(\sqrt x)\)、\(h/\sqrt x\to\lambda\in\mathbb R\)，以及
\(\sqrt x=o(|h|)\)、\(|h|=o(x^{2/3})\) 时的 \(R(x,h)\)。最后固定 \(c\in\mathbb R\setminus\{0\}\)，令 \(h=cx^{2/3}\)，计算
\[
 \lim_{x\to\infty}e^{-h^2/(2x)}R(x,h),
\]
说明上面的有效范围为什么不能直接扩大为 \(|h|=O(x^{2/3})\)。
\end{exercise}
\begin{solution}
写正实轴 Stirling 公式的对数形式
\[
 \log\Gamma(t)=\left(t-\frac12\right)\log t-t
                 +\frac12\log(2\pi)+r(t),\quad r(t)\to0.
\]
对 \(|h|\le L(x)\)，充分大的 \(x\) 满足 \(|h|\le x/2\)，于是 \(x+h\ge x/2>0\)。两次应用此式并约去共同项，得到
\[
 \log R(x,h)=\left(x+h-\frac12\right)
             \log\left(1+\frac hx\right)-h
             +r(x+h)-r(x).
\]
由于
\(\sup_{t\ge x/2}|r(t)|\to0\)，最后两项之差对整个 \(h\) 区间一致趋于零。这一步仅使用 Stirling 余项趋零，不对余项作未经证明的微分。

在 \(|u|\le1/2\) 上，Taylor 余项统一满足
\(\log(1+u)=u-u^2/2+O(|u|^3)\)。代入 \(u=h/x\)，整理得到
\[
 \left(x+h-\frac12\right)\log\left(1+\frac hx\right)-h
 =\frac{h^2}{2x}
   +O\left(\frac{|h|}x+\frac{|h|^3}{x^2}\right),
\]
其中常数不依赖于 \(x,h\)。所以所求上确界不超过
\[
 C\left(\frac{L(x)}x+\frac{L(x)^3}{x^2}\right)
 +2\sup_{t\ge x/2}|r(t)|,
\]
这在 \(L=o(x^{2/3})\) 时趋于零。

因而若 \(h=o(\sqrt x)\)，则 \(R(x,h)\to1\)；若 \(h/\sqrt x\to\lambda\)，则
\(R(x,h)\to e^{\lambda^2/2}\)。特别地，\(h=\sqrt x\) 时极限为 \(e^{1/2}\)。若
\(\sqrt x=o(|h|)\) 且 \(|h|=o(x^{2/3})\)，则 \(h^2/(2x)\to\infty\)，故 \(R(x,h)\to\infty\)。这些结论对正、负平移都适用，只要保持 \(x+h>0\)。

最后取 \(h=cx^{2/3}\)。这时必须把对数保留到三次：
\[
 \log(1+u)=u-\frac{u^2}2+\frac{u^3}3+O(u^4).
\]
代入同一恒等式，有
\[
 \log R(x,h)-\frac{h^2}{2x}
 =-\frac{h}{2x}-\frac{h^3}{6x^2}
   +O\left(\frac{h^2}{x^2}+\frac{|h|^4}{x^3}\right)+o(1)
 \longrightarrow-\frac{c^3}6.
\]
于是所求极限为 \(e^{-c^3/6}\ne1\)。在临界尺度，三次项保留了非零贡献；固定平移公式中可以忽略的项，在增长参数下必须重新估计。
\end{solution}
